\subsection{The precise parametric extent of the PL comparison}

\paragraph{Scope of the concordance computation.}
A.131 supplies whole effective countable simplicial models for the
PL and semialgebraic concordance spaces of every closed finite PL
model, and complete recursive presentations of all their homotopy
groups. Its component presentation also computes the stipulated
definable component group in every o-minimal expansion. The checked
PL--TOP comparison gives the same whole model for real TOP
concordances when every base component has dimension at least five,
and A.134 also supplies this model for every closed base whose
components have dimension at most two.
These outputs are exhaustive effective presentations; no finite
classical invariant decomposition or terminating word problem is
claimed. The earlier $E(K)$ presentation and the sphere calculation
supply additional structural information in their stated ranges.
A.132 refutes the unrestricted real PL--TOP weak-equivalence
claim: for $S^3$ the TOP concordance space has $\pi_2=\mathbb Z/2$,
whereas the jointly PL space has $\pi_2=0$. A.160 now proves that
the actual PL-to-definable comparison is a weak equivalence over
every fixed o-minimal real closed field, including the semialgebraic
case. A.161 combines that theorem with A.115 and A.131 to give the
whole effective definable-space model and complete recursive
presentations, with the component group first. The remaining A.1(b)
gate is the precise real TOP computation and comparison in base
dimensions three and four. A.142--A.146 give further component
obstructions and realizations there without evaluating the full
kernel or whole space. Earlier restricted criteria below remain
valid; their former comparison questions are resolved by A.160.

\begin{saresult}{A.114}{Proved}{The concordance comparison is surjective on the fundamental group}
Let $B$ be a closed compact PL manifold, and use total-space PL and
definable families to define the respective simplicial concordance
groups $P_{\rm PL}(B)\subset P_{\rm def}(B)$. In addition to the
bijection on components in A.113, inclusion induces a surjection
\[
 \pi_1(P_{\rm PL}(B),\operatorname{id})
 \longrightarrow
 \pi_1(P_{\rm def}(B),\operatorname{id}).
\]
More explicitly, every based definable circle family of concordances
is based definably homotopic to a family that is PL as a map of
total spaces.
If the field is $\R$ and every component of $B$ has dimension at
least five, this map on $\pi_1$ is an isomorphism. In that scope
the inclusion is injective on every homotopy group.
For $k\ge2$ there is the precise splitting
\[
 \pi_kP_{\rm def}(B)\simeq
 \pi_kP_{\rm PL}(B)\oplus
 \ker\bigl(\pi_kP_{\rm def}(B)\to\pi_kP_{\rm TOP}(B)\bigr).
\]
\end{saresult}
\begin{proof}
Write $Z=B\times[0,1]$ and let $C$ be a finite polygonal circle
with a distinguished vertex $c_0$. A loop is a homeomorphism
\[
 H:Z\times C\longrightarrow Z\times C
\]
over $C$, fixing $B\times\{0\}\times C$ pointwise and equal to
the identity over $c_0$. Mark also the upper boundary component.
Choose a PL function $a:C\to K$ nonconstant on every edge; assigning
distinct values to the finitely many vertices suffices. Apply the
complement to Shiota's Theorem 2.2 to the scalar function
$f=a\circ\operatorname{pr}_C$, with its two triangulations
$\pi=\operatorname{id}$ and $\pi'=H$. Mark the compact subpolyhedra
$Z\times e$ for every closed edge $e$ of $C$, and the two boundary
components of $Z$ times $C$. The source statement is
\cite[Theorem 2.2, Complement, p.~7]{Shiota2010}.

It supplies a definable isotopy $\omega_s$ from $H$ to a PL
homeomorphism, preserving $f$ and every marked set. Each $\omega_s$
preserves the full circle coordinate: for a point with coordinate
$c\in e$, its image has coordinate in $e$ with the same value of
$a$, and $a|_e$ is injective. This also works at vertices.
Thus $\omega_s$ is a family of homeomorphisms of $Z$ over $C$.
It need not yet fix the bottom pointwise. Write its bottom
restriction as
\[
 (b,0,c)\longmapsto(\alpha_s(b,c),0,c)
\]
and extend this restriction over the interval by
\[
 E_s(b,r,c)=(\alpha_s(b,c),r,c).
\]
This is a definable family of homeomorphisms. At $s=1$ it is
PL, since it is the product extension of the restriction of the
PL map $\omega_1$. The homeomorphisms
$g_s=E_s^{-1}\omega_s$ therefore fix the bottom and preserve the
top. At $s=0$ the correction is the identity, so $g_0=H$.
At $s=1$, $g_1$ is PL as a total-space map.

Finally let $q_s=g_s|_{Z\times\{c_0\}}$, viewed as a concordance
of $B$, and replace $g_s$ by
\[
 h_s=(q_s^{-1}\times\operatorname{id}_C)\circ g_s.
\]
Then $h_s$ is based at the identity
over $c_0$ for every $s$. Since $q_0=\operatorname{id}$, it starts
at $H$. Since $q_1$ is a restriction of a PL map and its product
extension is PL, it ends at a jointly PL family. These constructions
are continuous on total spaces and definable; inverse families
are continuous by compactness, or directly by the displayed
compositions. They prove the asserted based loop homotopy.

Loops and homotopies in these simplicial groups are represented,
after finite simplicial subdivision, by precisely such finite
parameter families. Thus the statement about families proves
surjectivity on $\pi_1$ of their realizations.

Over $\R$ a definable family is a continuous family and therefore
defines a map to $P_{\rm TOP}(B)$. The composite
\[
 P_{\rm PL}(B)\longrightarrow P_{\rm def}(B)
 \longrightarrow P_{\rm TOP}(B)
\]
is the usual PL--TOP comparison. Under the dimension hypothesis
it is a weak equivalence by A.113. Its maps on homotopy groups
are injective, so the first factor is injective on every homotopy
group. The surjectivity already proved gives the assertion on
$\pi_1$. This argument is confined to the real field, where
the displayed forgetful map to topological families exists.
If $i_k$ and $j_k$ denote the two maps on $\pi_k$ in that
display, $(j_ki_k)^{-1}j_k$ is a retraction of $i_k$.
For $k\ge2$ these groups are abelian, so its kernel gives the
stated direct-sum splitting. It is exactly $\ker j_k$.
\end{proof}

\begin{saresult}{A.115}{Proved}{All PL concordance homotopy groups are independent of the real closed field}
Let $F\subset K$ be an extension of real closed fields and let $B$
be a closed compact PL manifold given by a finite polyhedral model
over $F$. Extension of coefficients induces a weak equivalence
of simplicial groups
\[
 P_{{\rm PL},F}(B)\longrightarrow P_{{\rm PL},K}(B_K).
\]
Consequently a fixed finite abstract PL model has the same PL
concordance homotopy type over any two real closed fields, by the
zigzag through a realization over the field of real algebraic
numbers. This is a zigzag of coefficient-extension maps; it is
not a literal map from a non-Archimedean field's families to
ordinary topological families.
\end{saresult}
\begin{proof}
We give the parameter-preserving argument. A statement only about
unparametrized homeomorphisms would not suffice.

\emph{Extension on product cells.}
Let $Z$ be the underlying space of a finite complex of convex
cells, let $Q$ be a finite parameter polyhedron and $Q_0$ a
subcomplex. Suppose a PL homeomorphism of $Z\times Q_0$ is over
$Q_0$ and preserves each cell of $Z$. It extends to a PL
homeomorphism of $Z\times Q$ over $Q$, still preserving every
cell. More generally it extends relative to any already
specified union of product cells on which the same conditions hold.
To see this, triangulate $Q$ and proceed over the cells
$\sigma\times\tau$ in increasing total dimension. On the boundary
of a cell not already specified the homeomorphism has been
defined, maps that boundary to itself, and preserves the parameter
coordinate. Choose an interior point
$c=(c_\sigma,c_\tau)$ of this convex product cell and extend
by coning:
\[
 c+t(z-c)\longmapsto c+t(H(z)-c),
 \qquad z\in\partial(\sigma\times\tau),\quad 0\le t\le1.
\]
This is PL after starring a triangulation of the boundary at
$c$, and is a homeomorphism by the same construction for the
inverse boundary map. Its parameter coordinate is unchanged,
since
\[
 \operatorname{pr}_Q(c+t(H(z)-c))
 =(1-t)c_\tau+t\operatorname{pr}_Q(z)
 =\operatorname{pr}_Q(c+t(z-c)).
\]
Thus it is coning on the \emph{total product cell}, not separate
radial coning in each fibre. The maps agree on faces and the
finite induction proves the claim. It works over every ordered
field. In particular cell-preserving PL families can be contracted
through such families. Two families agreeing on a union of
product cells can be contracted with that agreement retained:
first perform this same extension on the common part, then on
the remaining cells separately.

Put $Z=B\times[0,1]$, with a fixed finite convex-cell
decomposition $\mathcal C_Z$ over $F$ that contains its bottom
and top as subcomplexes. Represent a homotopy group element by
a PL homeomorphism
\[
 h:Z_K\times P_K\longrightarrow Z_K\times P_K,
 \qquad P=[0,1]^q,
\]
over $P_K$, fixing the bottom and equal to the identity on
$Z_K\times\partial P_K$. For $q=0$ the last condition is empty.
Fix a cell decomposition $\mathcal C_P$ of $P$ over $F$.

\emph{Finite transferred data.}
Apply simultaneous subdivision to the one-way chain consisting
of two distinct copies of the total space,
\[
 Z_K\times P_K\xrightarrow{h}Z_K\times P_K
          \xrightarrow{\operatorname{pr}_P}P_K.
\]
The source and target subdivisions $M,N$ can be required to
refine the same fixed product cells
$\mathcal C_Z\times\mathcal C_P$, and the base subdivision $L$
refines $\mathcal C_P$. Then $h:M\to N$ is a simplicial
isomorphism and both projections to $L$ are simplicial, the
source projection being the composite of the displayed maps.
This is exactly the permitted one-way case of
\cite[Theorems 2.14--2.15, pp.~17--19]{RourkeSanderson1972};
it does not ask for a simplicial subdivision of a cycle.
On the bottom and on $Z_K\times\partial P_K$, where $h$ is
the identity, the induced source and target subdivisions coincide.

Keep the finite abstract complexes and all their carrier labels
fixed, but regard their vertex coordinates as unknowns. The
requirements that they realize subdivisions of the designated
fixed product cells, that the abstract map $M\to N$ is realized
by a PL homeomorphism, and that each source vertex and its
corresponding target vertex project to their designated vertex
of $L$, are a finite first-order condition over $F$. To spell
out the finiteness, affine independence is expressed by nonzero
minors; simplex membership uses finitely many barycentric
coordinates; intersection along the prescribed common faces
and coverage of each fixed cell are finite quantified field
formulas. Add equality of the corresponding source and target
vertices on the bottom and on the parameter boundary. The
given configuration satisfies this condition over $K$.
Quantifier elimination for real closed fields, hence
$F\preccurlyeq K$, gives a configuration $M',N',L'$ over $F$
satisfying the same condition. It determines a PL family
$h_F$ of concordances, based on $\partial P$, over $F$.

There are matching-vertex PL homeomorphisms over $K$
\[
 u:M\longrightarrow M',\qquad
 v:N\longrightarrow N',\qquad
 \beta:L\longrightarrow L',
 \qquad h_F=vhu^{-1}.
\]
Each of $u,v$ preserves every fixed product cell, and $\beta$
preserves every fixed base cell. The projection identities
hold on entire simplices, not just their vertices:
\[
 \operatorname{pr}_P u=\beta\operatorname{pr}_P,\qquad
 \operatorname{pr}_P v=\beta\operatorname{pr}_P.
\]
Indeed a simplex of $M$ projects into one simplex of $L$,
on which $\beta$ is affine. For
$x=\sum_i\lambda_i x_i$ in this simplex, the vertex-incidence
equalities give
$\operatorname{pr}_P u(x)=\sum_i\lambda_i\beta(\operatorname{pr}_P x_i)
=\beta(\operatorname{pr}_P x)$; the proof for $N$ is identical.

Set
\[
 u_0=(\operatorname{id}_Z\times\beta^{-1})u,\qquad
 v_0=(\operatorname{id}_Z\times\beta^{-1})v.
\]
They preserve the parameter pointwise and every product cell.
They agree on the bottom and on $Z_K\times\partial P_K$,
because the matching-vertex maps agree there. Apply the
product-cell extension construction over the enlarged parameter
$P_K\times[0,1]$ to obtain PL isotopies $u_s,v_s$ from the
identity to $u_0,v_0$ which retain this agreement for every $s$.
Concretely extend their common boundary data first on the union
of these product subcomplexes, then extend over the other cells.
The families
\[
 v_s h u_s^{-1}
\]
are therefore concordances fixed on the bottom and based on
the parameter boundary throughout. Their last member is
$(\operatorname{id}_Z\times\beta^{-1})h_F
(\operatorname{id}_Z\times\beta)$.
The same product-cell construction gives an isotopy of the
cell-preserving map $\beta$ to the identity, preserving every
base cell. Conjugating by that isotopy connects this last member
to $h_F$ while keeping the parameter boundary based. This proves
surjectivity on every homotopy group, including components.

\emph{Injectivity.}
Suppose a based PL family over $F$ becomes nullhomotopic over
$K$. Take one finite PL triangulation of its $K$-nullhomotopy.
Its vertex coordinates again give a finite first-order
configuration problem over $F$, with the extra requirement that
on the parameter boundary its PL map equal the prescribed
$F$-family or the identity, as appropriate. That equality is
itself a finite field formula, because the prescribed family
has finitely many affine pieces over $F$. Transfer therefore
gives an $F$-PL nullhomotopy with exactly the same boundary data.
The same argument applies to a path between two zero-simplices.
Thus the maps on all homotopy groups are injective.

Finally simplicial groups are Kan, and their homotopy groups
are computed by these finite parameter families. The established
isomorphisms prove the weak equivalence. Real algebraic numbers
embed as an ordered real closed subfield in every real closed
field; realizing a finite abstract complex by standard rational
vertices gives the asserted comparison zigzag.
\end{proof}

\begin{saresult}{A.116}{Proved}{Semialgebraic concordance comparison reduces to the real field}
Let $F\subset K$ be real closed fields and let $B$ be a closed
compact PL manifold given by a finite polyhedral model over $F$.
Extension of coefficients induces a weak equivalence
\[
 P_{{\rm SA},F}(B)\longrightarrow P_{{\rm SA},K}(B_K),
\]
where families and their homotopies are semialgebraic maps of
total spaces. The squares formed by these maps and the PL
coefficient-extension maps commute. Thus the homotopy-group
comparison $P_{\rm PL}\to P_{\rm SA}$ for a finite abstract PL
model is independent of the real closed field and reduces to $\R$.

In particular, if every component of $B$ has dimension at least
five, then over every real closed field this comparison is
bijective on components, is an isomorphism on $\pi_1$, and is
injective on every homotopy group. Full weak equivalence is not
asserted here.
\end{saresult}
\begin{proof}
We give the continuity argument needed in addition to transfer.
Merely transferring the existence of a homeomorphism would not
connect its coefficient tuple to the original family.

\emph{Continuous evaluation after a finite coefficient partition.}
Let $D$ be a closed bounded semialgebraic set over $F$, and fix
one ordered-field formula $\varphi(a;x,y)$ over $F$. Let $T$
consist of the tuples $a$ for which this formula is the graph of
a homeomorphism $h_a:D\to D$ satisfying a prescribed finite
list of first-order conditions. Totality, uniqueness,
bijectivity, and continuity of both the map and its inverse are
first-order conditions: continuity uses quantifiers over points
and positive $\epsilon,\delta$ in the field. Quantifier
elimination therefore makes $T$ semialgebraic over $F$.

Form the graph space
\[
 \Gamma=\{(a,x,y):a\in T,\ x,y\in D,\ \varphi(a;x,y)\}.
\]
Its projection to $T$ is continuous and semialgebraic, although
evaluation on $T\times D$ need not yet be continuous. Apply
semialgebraic triviality to this projection over $F$, and refine
the resulting base partition into semialgebraic cells over $F$.
The precise input is Hardt triviality without a properness
assumption on the original map; its real-closed-field form is
stated in \cite[Section 2.2.1, p.~1296]{LerarioRizziTiberio2025}.
For each cell $C$, choose $a_0\in C(F)$ and identify the fixed
fibre of its trivialization with the fibre at $a_0$. We obtain
\[
 \theta:C\times\Gamma_{a_0}\xrightarrow{\ \sim\ }\Gamma|_C
\]
over $C$. Choose all these data over $F$ before extending to
$K$. Their finite first-order properties, including continuity,
remain true after coefficient extension.

The fibre $\Gamma_{a_0}$ is closed and bounded, being the graph
of a continuous map on $D$. Compose $\theta$ with the projection
which forgets the output coordinate to obtain a continuous
semialgebraic bijection
\[
 J:C\times\Gamma_{a_0}\longrightarrow C\times D.
\]
It is a homeomorphism. Indeed, for a closed bounded
semialgebraic $A\subset C\times D$, its projection
$C_0\subset C$ is closed and bounded, and $J^{-1}(A)$ is
closed in $C_0\times\Gamma_{a_0}$, hence closed and bounded.
Thus $J$ is definably proper. Inverse continuity can also be
checked directly: curve selection applied to a failure of
continuity would give a semialgebraic curve tending to a point
of $C\times D$ whose inverse fails to tend to the unique
inverse point. Its second inverse coordinate lies in
$\Gamma_{a_0}$ and has a limit by boundedness and semialgebraic
monotonicity. Continuity of $J$ identifies that limit with the
unique inverse point, a contradiction. This uses definable
compactness, not ordinary compactness of the order topology
of a non-Archimedean field.

The output projection composed with $\theta J^{-1}$ is
$(a,x)\mapsto h_a(x)$. Thus evaluation is jointly continuous on
$C\times D$. Interchanging the input and output coordinates
proves the same assertion for inverse evaluation. We have
obtained genuine families of homeomorphisms on the coefficient
cells.

\emph{Surjectivity on homotopy groups.}
Put $Z=B\times[0,1]$ and $D=Z\times P$, where $P=[0,1]^q$.
Represent a homotopy-group class over $K$ by a semialgebraic
homeomorphism $h:D_K\to D_K$ over $P_K$, fixed on the bottom
and equal to the identity over $\partial P_K$. For $q=0$ the
boundary condition is empty. Choose one finite formula for its
graph and replace its finitely many extra coefficients by a
tuple $a$, leaving $D$ and the formula fixed over $F$. In $T$
impose exactly the conditions just stated: preservation of each
parameter point, the bottom fixed pointwise, the top preserved,
and the identity on the parameter boundary.

The original tuple lies in $T(K)$, hence in $C(K)$ for one of
the cells chosen over $F$. Choose $a_0\in C(F)$. A semialgebraic
cell has a semialgebraic homeomorphism over $F$ with an open
cube (or is a point). Transporting a line segment in that cube
gives a semialgebraic path in $C(K)$ from the original tuple to
$a_0$. Pull back the jointly continuous evaluations along this
path. They give a semialgebraic homotopy of concordance families
through members satisfying all conditions encoded in $T$.
The last family is defined over $F$. This proves surjectivity
on every homotopy group and on components.

\emph{Injectivity on homotopy groups.}
Suppose a based family over $F$ admits a semialgebraic
nullhomotopy over $K$. Choose one finite graph formula for that
nullhomotopy and replace its extra coefficients by variables.
Require a total continuous family of homeomorphisms with
continuous inverse, preserved parameters, fixed bottom and
parameter sides, and exact equality with the prescribed
$F$-family and the identity at the two endpoints. These are
first-order conditions over $F$, including equality with the
prescribed endpoint graph. A solution exists in $K$, so
$F\preccurlyeq K$ supplies a solution in $F$. This is the
required $F$-semialgebraic nullhomotopy, with its boundary
unchanged. The same argument descends paths between two
zero-simplices. Thus every homotopy-group map is injective.
As in A.115, the Kan property and finite parameter
representatives give weak equivalence of simplicial groups.

Coefficient extension preserves PL maps, so the PL/SA squares
commute literally. Realize the finite abstract PL model over
the real algebraic numbers and use their embeddings in $K$
and $\R$. A.115 and the present theorem give vertical weak
equivalences in the comparison zigzag. Over $\R$, A.114 proves
injectivity in every degree and surjectivity on $\pi_1$ in the
stated dimension range; A.113 supplies the component assertion.
These assertions therefore hold over $K$. This reduction is
for semialgebraic structures; it does not assume that an
arbitrary o-minimal expansion has a model whose field is $\R$.
\end{proof}

\paragraph{Exact remaining gate; no full parametric conclusion.}
The scalar straightening proof in A.114 uses the fact that a scalar function is injective on
each edge of a finite graph. It does not generalize in this way
to a two-dimensional parameter disk. For an arbitrary o-minimal
expansion of a real closed field it does not prove injectivity on $\pi_1$, and it does not prove that
$P_{\rm PL}(B)\to P_{\rm def}(B)$ is a weak equivalence.
A first missing assertion is this relative extension statement:
given a definable family on $D^2$ whose boundary family is jointly
PL, can it be deformed relative to its boundary to a jointly PL
family while preserving every point of $D^2$ as a parameter?
Over $\R$ in the dimension range of A.113, the additional TOP
argument above already settles injectivity. By A.115--A.116 this
also holds for semialgebraic families over every real closed field.
In this semialgebraic scope the remaining comparison reduces to
the real field, and the first remaining group assertion is
surjectivity on $\pi_2$. This requires
straightening a sphere family with two independent parameters,
or another argument that rules out a definable class killed by
the forgetful map to TOP.

Triangulating a total-space homeomorphism with the finitely many
parameter faces marked only preserves those faces setwise. It
does not preserve every fibre of the projection to $D^2$.
Successively applying the scalar theorem to its two coordinates
does not repair the problem, because the second isotopy need not
preserve the first coordinate. General homotopy comparison for
maps of definable sets likewise does not keep intermediate maps
homeomorphisms. These are exact failures of proposed deductions,
not counterexamples to the desired weak equivalence. The obstruction
to triangulating arbitrary maps with a two-dimensional target is
explicitly discussed in \cite[p.~1231]{FichouShiota2017}; the weak
triangulation used there is local and does not supply the missing
parameter-preserving family theorem.

The local-contractibility arguments used for A.117 and A.209 take
place in TOP homeomorphism groups. Their continuous contractions
and quotient lifts are not asserted semialgebraic. In particular,
using such a contraction to fill a small semialgebraic sphere would
require a separate proof that the filling is jointly semialgebraic.
Individual Hauptvermutung, cellwise Hardt triviality, or a collection
of unrelated PL approximations does not establish this. The
compact-bundle reduction suggested by clutching a sphere family
has the same missing parameter-preserving reduction, so it is not
an independent proof.

The paragraph after Theorem 2.6 in
\cite[p.~35]{Shiota2010} proposes a PL reduction for definable
fibre bundles as a further possibility. Its Remark 2.7 proves
only stable equivalence for microbundles. Neither statement is
the missing relative two-parameter theorem. Similarly, the positive
bundle triangulation statement in \cite[p.~160]{MurayamaShiota2013}
assumes a genuine $C^k$ fibre bundle, $k>0$; the local-to-global
question discussed on pp.~161--162 is separate.
The finite-atlas semialgebraic bundle results in
\cite[Definition~8.1 and Propositions~8.4--8.5, Lemmas~8.7--8.8,
pp.~2529--2532]{HardtLambrechtsTurchinVolic2011} prove pullback,
composition and cylinder invariance; they do not compare these bundles
with PL bundles. Likewise, the uniform bounds in
\cite[Proposition~3.6 and Theorem~3.8, pp.~1582--1585;
Theorem~7.3, pp.~1595--1596]{DemdahKartoue2011} concern individual
semialgebraic homeomorphisms. Their proof uses Hardt pieces and finitely
many chosen representative maps. It provides no continuous gluing of
those choices across the boundaries of the pieces. Such gluing over a
two-dimensional parameter disk is precisely an additional requirement
here.

There is also a concrete obstruction to applying the finite
fragmentation construction of \cite[Proposition~52, pp.~84--85;
Lemma~12, pp.~90--93]{Mnev2007} verbatim to semialgebraic input.
Proposition~52 asserts that sufficiently small affine-on-simplex
Hudson graphs meet each leaf of a jointly PL family once. Let
$X=[-1,1]$, $-1/8\le t\le1/8$, and
\[
 h_t(x)=x+\sqrt[3]{t}(1-|x|).
\]
This is a jointly continuous semialgebraic family of homeomorphisms
fixing both endpoints: its two slopes are $1+\sqrt[3]{t}$ and
$1-\sqrt[3]{t}$, both at least $1/2$. The leaf through $x=0$ is
$x=\sqrt[3]{t}$. For every $0<s<1/8$, the affine graph $t=sx$
meets it three times, since $x^3=sx$ has the three distinct
solutions $0,\sqrt{s},-\sqrt{s}$. Replacing $t$ by
$\tau=1/2+4t$ puts the base in $[0,1]$; the corresponding Hudson
function $\tau(x)=1/2+4sx$ has diameter $8s$, which is
arbitrarily small. In particular $s=1/16$ gives the exact rational
intersections $(x,t)=(0,0),(1/4,1/64),(-1/4,-1/64)$, checked in
\texttt{check\_hudson\_failure.py}. Thus this PL lemma cannot
simply be relabelled semialgebraic. Nonlinear graphs adapted to the
family might still provide a different proof; none has been
constructed here. This counterexample refutes that proposed
extension of a lemma, not the desired concordance comparison.

Even making all parameter tracks locally Lipschitz cannot in general
be achieved by a semialgebraic change of the base coordinate alone.
For a precise real-field example let
\[
 c(x)=\frac{x}{1+|x|},\qquad
 g_t=c\circ\bigl(x\longmapsto\sqrt[3]{x^3+t}\bigr)\circ c^{-1}
       \quad\hbox{on }(-1,1),
\]
and extend $g_t$ by fixing the endpoints. For $t$ in any compact
interval this is a jointly continuous semialgebraic family of
homeomorphisms of $[-1,1]$. If $t=\lambda(s)$ is any increasing
semialgebraic homeomorphism of nondegenerate intervals, it has an
interior point $s_0$ where $0<\lambda'(s_0)<\infty$.
Choose $x_0^3=-\lambda(s_0)$. The track at $c(x_0)$ then equals
\[
 g_{\lambda(s)}(c(x_0))
      =c\bigl(\sqrt[3]{\lambda(s)-\lambda(s_0)}\bigr).
\]
Its absolute value divided by $|s-s_0|$ tends to infinity, so it is
not locally Lipschitz at $s_0$. Decreasing reparameterizations give
the same conclusion with the sign reversed. This only rules out
that base-only regularization method; a change depending on the
spatial variable, or a different adapted graph construction, has
not been excluded.

The sources opened for this check, together with the arguments above, justify
A.113--A.116, but do not
settle the full higher-homotopy comparison.

\subsection{Pointwise cone families with locally parameterized outer faces}

There is a useful intermediate family convention between a derived
choice space and Hausdorff continuity of subsets alone. Here an outer
face is an unparameterized subset. A finite-parameter family is required
locally on its parameter space to admit jointly continuous embeddings
of a fixed abstract outer face onto those subsets. The parameterizations
are witnesses of continuity, not additional data retained in the objects.
In a definable simplicial family model the witnesses on each simplex
are jointly definable; finite nullhomotopies use finitely many such
simplex presentations. No finite subcover assertion over a
non-Archimedean field is implicit in this convention.

\begin{saresult}{A.117}{Refuted}{The cone obstruction persists for locally parameterized outer-face families}
Let $B$ be a closed connected PL manifold of dimension at least four,
let $X=C_2B$ be its cone with radial interval $[0,2]$, and put $C=C_1B$.
Give the locally flat outer faces of strict cone neighborhoods the
local quotient topology from
\[
 \operatorname{Emb}_{\rm lf}\bigl(B,B\times(0,2)\bigr)
       /\Homeo(B),
\]
where embeddings have the compact-open topology and $\Homeo(B)$ acts
by precomposition. An outer face determines the side containing the
vertex. The component $\mathcal W$ of $C$ has weak homotopy type
\[
             \mathcal W\simeq BP_{\rm TOP}(B).
\]
For $B=(S^1)^6$, it is therefore not contractible. Its nontrivial loop
can be chosen semialgebraic, and it cannot be filled by real definable
families of strict cones whose outer faces admit local joint
parameterizations. The same non-filling conclusion holds for finite
semialgebraic simplex families with these joint semialgebraic witnesses
over every real closed field.

If an external parameterization $B\to\dot C$ is retained as part of
each object, the component instead has weak homotopy type
\[
 B\Homeo\bigl(B\times[1,2]\ \mathrm{rel}\
               B\times\{1,2\}\bigr).
\]
This is a different convention from the unparameterized outer faces
in the preceding assertion.
\end{saresult}
\begin{proof}
We first record the topological lifting fact used below. If $G$ is a
metrizable topological group and $J$ is a closed, completely metrizable
subgroup that is locally $n$-connected for every finite $n$, every
continuous map $[0,1]^m\to G/J$ lifts to $G$. This is the finite-cube
lifting argument in the proof of
\cite[Corollary~2.2, pp.~4--5]{BanakhRepovs2010}. The additional
``balanced'' hypothesis in that corollary is used for its separate
homogeneous-metric conclusion, not for this lifting argument. Its
hypotheses can be checked directly: the cosets are closed complete
subsets in the two-sided completion, and the coset family is equi-locally
$n$-connected. For the latter, given an identity neighborhood $U$, choose
$W$ so that every sphere in $J\cap W$ contracts in $J\cap U$, and then
choose $V$ with $V^{-1}V\subset W$. Translating a sphere contained in
$gJ\cap xV$ by the inverse of any one of its points reduces it to that
contraction in $J$. The finite-dimensional selection theorem used in
the cited proof supplies lifts over cubes. In particular the assertion
uses complete metrizability, not completeness of a one-sided invariant
metric. A lift of a homotopy of cubes can be made to extend a prescribed
initial lift: multiply it on the right by its initial discrepancy in
$J$, constant in the homotopy coordinate. Thus $G\to G/J$ is a Serre
fibration.

Now set
\[
 G=\Homeo(X\ \mathrm{rel}\ \{o\}\cup B\times\{2\}),
 \qquad J=\{g\in G:g(C)=C\}.
\]
The Alexander contraction is a continuous strong contraction of $G$.
Restriction to the shell $B\times[1,2]$ gives a split map
$J\to P_{\rm TOP}(B)$, with section obtained by coning the inner-face
restriction. Its kernel is
\[
 K=\Homeo(C\ \mathrm{rel}\ \{o\}\cup B\times\{1\}).
\]
Restriction, conical extension, and composition are continuous in the
compact-open topology, including inverses. Hence this splitting is a
homeomorphism of underlying spaces $J\cong K\times P_{\rm TOP}(B)$.
The factor $K$ has the continuous strong Alexander contraction.
Such a contraction proves local contractibility at the identity:
for a specified identity neighborhood, continuity and compactness of
the contraction interval provide a smaller neighborhood whose entire
contraction stays in the specified one. Translation gives the assertion
at every point.

The other factor is locally contractible as well. The boundary-component
preserving subgroup of $\Homeo(B\times[1,2])$ is open and is locally
contractible by \cite[Corollary~1.1, p.~63]{EdwardsKirby1971}.
Restriction to the top boundary has the continuous section given by
constant extension across the interval. The resulting continuous
retraction onto the top-fixed subgroup shows that $P_{\rm TOP}(B)$
is locally contractible. Finally the metric
\[
 d_G(g,h)=\sup_{x\in X}d(gx,hx)
           +\sup_{x\in X}d(g^{-1}x,h^{-1}x)
\]
is complete for any complete compatible metric on the compact $X$.
A Cauchy sequence has mutually inverse uniform limits. The prescribed
fixed-point conditions and the equation $g(C)=C$ are closed in this
metric. Thus $G$ and $J$ meet the lifting hypotheses.

We next verify that the orbit topology is the stipulated outer-face
topology. The codimension-one local unknotting theorem
\cite[Theorem~7.3.1, p.~358]{DavermanVenema2009} says that sufficiently
close locally flat embeddings of the compact face are related by an
arbitrarily small ambient push. Here closeness is uniform closeness of
\emph{parameterized embeddings}. An $\epsilon$-push is supported in
the $\epsilon$-neighborhood of the face
\cite[p.~115]{DavermanVenema2009}. Choose $\epsilon$ smaller than its
distance from both $o$ and $B\times\{2\}$. The push then extends by the
identity to $X$ and belongs to $G$; its inverse is equally small. It
sends the cone side to the cone side because it fixes $o$.

More explicitly, let $U\subset G$ be open and $g\in U$. For every
parameterization $e:B\to g(\dot C)$, a sufficiently close embedding
$e'$ is $he$ for a push $h$ sufficiently close to the identity that
$hg\in U$. Therefore the inverse image of the orbit image of $U$ in
the parameterized embedding space is open. Passing to the quotient by
precomposition proves that $G\to G\cdot C$ is open. It consequently
induces a homeomorphism $G/J\cong G\cdot C$. The same argument at any
strict cone with a locally flat outer face shows that every ambient
orbit is open. As $G$ is contractible, its orbit is connected and is
exactly the component $\mathcal W$ under discussion.

The Serre fibration $J\to G\to\mathcal W$, with contractible middle
space, gives the asserted weak homotopy type and, in particular,
\[
 \pi_{k+1}(\mathcal W,C)\cong\pi_kJ
          \cong\pi_kP_{\rm TOP}(B)\qquad(k\geq0).
\]
Equivalently, the map from the homotopy quotient of the contractible
$G$ by $J$ to $G/J$ is a weak equivalence, identifying it with $BJ$.
No local normal charts preserving all family parameters were assumed.

For the explicit obstruction take a PL concordance $p$ of the six-torus
representing a nonzero class in $\pi_0P_{\rm TOP}(B)$; such a class and
its PL representative are supplied by A.103 and A.113. Regard $p$ as
an automorphism of the shell fixed on its top, and cone its inner-face
map to obtain $j\in J$. Let $H_s$, $0\leq s\leq1$, be the Alexander
isotopy from the identity to $j$. Then
\[
                  C_s=H_s(C)
\]
is a semialgebraic loop of strict cones, since $C_0=C_1=C$.
It has jointly continuous face parameterizations on the parameter
interval; the two endpoint parameterizations need not coincide,
because they are not part of the unparameterized objects.

Suppose this loop bounded a disk of the indicated real definable
families. Each individual definable outer face is TOP locally flat:
compatible triangulation makes the ambient annulus and the face
polyhedra, the definable Hauptvermutung makes both PL manifolds, and
a codimension-one PL embedding is TOP locally flat. The last statement
follows by induction on links and the generalized Schoenflies theorem,
as in \cite[Theorem~1.2.1, p.~29]{DavermanVenema2009}.
Local joint face parameterizations make the disk map continuous in the
embedding-quotient topology. Its image lies in $\mathcal W$, because
ambient orbits are open. By the finite-cube lifting fact the disk map
has a continuous lift $L$ to $G$. On its boundary compare $L_s$ with
$H_s$. The path $a_s=L_s^{-1}H_s$ lies in $J$, and $L_0=L_1\in J$.
Thus $a_s$ joins $L_0^{-1}$ to $L_0^{-1}j$, so multiplication by $L_0$
gives a path from the identity to $j$ in $J$. Shell restriction would
make $p$ trivial in $\pi_0P_{\rm TOP}(B)$, a contradiction. A disk
assembled from finitely many such parameterized simplices gives the
same continuous disk map, so the argument also excludes nullhomotopies
in the associated simplicial family model.

For semialgebraic families over an arbitrary real closed field, choose
the torus and $p$ over the real algebraic numbers, using A.115 and the
nonzero real PL class just used. If a finite semialgebraic filling
existed over another real closed field, write its finitely many graph
formulas and coefficient tuples. Joint continuity, fibrewise embedding
of each parameterized face, avoidance of the vertex and outer base,
agreement of face images on shared parameter faces, and the prescribed
boundary images are all first-order conditions. Transfer gives the
same finite face family over the real algebraic numbers and then over
$\mathbb R$. The transferred individual faces are semialgebraic and
therefore TOP locally flat. Conicality of the transferred interiors is
not needed: all these faces lie in the same open ambient orbit as the
boundary loop. This contradicts the real argument. No assertion about
real models of arbitrary o-minimal expansions is used.

Finally, retaining an external parameterization replaces $J$ by the
subgroup $J^{\rm pt}$ fixing $\dot C$ pointwise. Its shell restriction
has kernel $K$ and target
$Q=\Homeo(B\times[1,2]\ \mathrm{rel}\ B\times\{1,2\})$.
The group $Q$ is locally contractible: double the shell along its
boundary, apply \cite[Corollary~7.3, p.~83]{EdwardsKirby1971} relative
to the two doubling faces, and restrict the resulting open
component-preserving subgroup, which is $Q\times Q$, to one half.
All metric and orbit arguments above apply with $J^{\rm pt}$ and
parameterized embeddings. The component is therefore weakly equivalent
to $BQ$. This verifies exactly why that additional convention changes
the stabilizer.
\end{proof}

\paragraph{The distinct Hausdorff-only convention.}
A.117 settles the comparison for outer-face families locally continuous
in the compact-open embedding quotient. A.118 supplies an explicit
semialgebraic family of polygonal disks proving that Hausdorff
continuity of both cones and outer faces does not imply local joint
face parameterizations. The two conventions are therefore genuinely
different. Whether the six-torus loop becomes fillable in the weaker
Hausdorff-only space is a separate comparison question; it is not
needed for the refutation of II b) under the simplicial-category
convention stipulated in the prompt, already proved in A.111.
Likewise A.117 supplies TOP lifts, not definable lifts, and therefore
does not settle the higher PL-to-definable concordance comparison in
A.114--A.116.


\subsection{A complete countable presentation for components in a stable range}

The following computation specifies an abelian group by generators and
all its relations. It does not assert finite generation, a finite
invariant-factor decomposition, or an algorithm deciding equality of
arbitrary words. In particular it supplies the extension law that is
not determined by the outer terms of the Hatcher--Wagoner--Igusa exact
sequence.

\begin{saresult}{A.119}{Proved under stated hypotheses}{An explicit countable component presentation}
Let $B$ be a closed definable manifold whose finite PL model $K$ is
smoothable and has dimension at least $9$ on every component. The
construction below, applied to each connected component of $K$, gives
a countable abelian presentation $E(K)$ and an isomorphism
\[
                   \pi_0P_{\rm def}(B)\cong E(K).
\]
The group operation on the right is addition of integral cycles. The
presentation depends only on the finite simplicial homotopy model $K$,
and is valid over every real closed field and every fixed o-minimal
expansion. It also gives the stable smooth concordance component group
for any compact smooth manifold, without a dimension restriction.
It makes no unstable assertion for a nonsmoothable PL manifold or for
an arbitrary manifold of dimension at most $8$.
\end{saresult}

\paragraph{The presentation.}
We first give the construction for connected $K$. Collapse a maximal
tree in its $1$-skeleton and use the resulting reduced finite simplicial
set. Let $G$ be its Kan loop group and put $R=\Sigma^\infty_+G$, the
spherical group ring. A countable simplicial model for these objects
and all the operations used below is specified in the proof.
Let
\[
 H_n=\operatorname{GL}_n(R),\qquad
 W_n=G^n\rtimes\Sigma_n.
\]
Here $H_n$ is the derived simplicial monoid of homotopy automorphisms
of the free right $R$-module $R^{\vee n}$; it is not the discrete
group $\operatorname{GL}_n(\Z[\pi_1K])$. Permutations of summands and
left multiplication by elements of $G$ define $W_n\to H_n$.
Stabilize by adjoining an identity summand and form
\[
       C=\operatorname*{hocolim}_n BH_n,\qquad
       D=\operatorname*{hocolim}_n BW_n,
       \qquad f:D\longrightarrow C.                 \tag{A.119.1}
\]
We use countable based simplicial telescope models. If necessary,
collapse a maximal tree to give each model a single vertex; the map
$f$ can equivalently be retained before this collapse, using the
based lift described below. All chains in what follows are normalized
simplicial chains with integral coefficients and finite support.

For a connected one-vertex simplicial set $Y$, put
\[
 L_Y=C_1(Y)/\partial C_2(Y)=H_1(Y;\Z).
\]
Its universal abelian covering $Y^{\rm ab}$ has $q$-simplices
$(a,\sigma)$ with $a\in L_Y$ and $\sigma\in Y_q$. The first face is
\[
 d_0(a,\sigma)=
 (a+[\sigma|[0,1]],d_0\sigma),
\]
and the other faces and all degeneracies retain $a$. For $q=0$ there
are no face maps. The triangle relations in $L_Y$ verify the
simplicial identities. For a model with several vertices the same
formula uses the edge labels obtained by joining the two endpoints
to the base vertex in the chosen maximal tree. This specifies a
connected cover corresponding to
$\ker(\pi_1Y\to H_1Y)$. The based map $f$ induces a based map
$f^{\rm ab}:D^{\rm ab}\to C^{\rm ab}$.

Write $Z_2(Y)=\ker(\partial:C_2(Y)\to C_1(Y))$. The generators of
$E(K)$ are symbols $[z]$ for every $z\in Z_2(C^{\rm ab})$. Its
complete list of relations is
\[
\begin{aligned}
 [z+z']&=[z]+[z'] &&(z,z'\in Z_2(C^{\rm ab})),\\
 [\partial c]&=0 &&(c\in C_3(C^{\rm ab})),\\
 [(f^{\rm ab})_*w]&=0 &&(w\in Z_2(D^{\rm ab})).
\end{aligned}                                                   \tag{A.119.2}
\]
Thus it is also the explicit quotient
\[
 E(K)=\frac{Z_2(C^{\rm ab})}
 {\partial C_3(C^{\rm ab})+(f^{\rm ab})_*Z_2(D^{\rm ab})}.
                                                               \tag{A.119.3}
\]
All faces in these formulas are ordinary bar, simplicial, and
telescope faces; in particular every boundary is the alternating sum
$\partial\sigma=\sum_i(-1)^id_i\sigma$. For a disconnected $K$
take the finite product of the presentations for its components.

\begin{proof}
\emph{Countable models determined by finite data.}
For the reduced simplicial set $X$ obtained from $K$, the Kan loop
group is explicitly
\[
 G(X)_q=F(X_{q+1})/\langle[s_0x]=1:x\in X_q\rangle,
\]
with the following homomorphisms on generators:
\[
 d_0[x]=[d_1x][d_0x]^{-1},\qquad
 d_i[x]=[d_{i+1}x]\ (i>0),\qquad
 s_i[x]=[s_{i+1}x].
\]
These are the convention of \cite[V.5, (5.1), p.~276]
{GoerssJardine1999}; their Theorem~5.10 and Corollary~5.11
(p.~285) give the loop-space comparison. All elements are finite
words in the given simplices, so the simplicial sets are countable.
Use simplicial symmetric spectra for $R$ and their usual
stable module model structure. Here are precise choices that avoid
using an uncountable singular mapping space. Let $\mathcal F_R$
be the category of finite cell right $R$-modules, with their cell
presentations retained. Cells are the free modules on
$F_j\partial\Delta^q_+\to F_j\Delta^q_+$. All attaching maps and all
maps between such objects are specified by the images of finitely
many cell generators in countable simplicial sets. Consequently the
category has a countable skeleton and countable morphism sets.

For acyclic attachments use the countable set $R\wedge J$, where
$J=FI_\Lambda\cup\{c_j\mathbin\square i\}$ is the explicit set of
\cite[Definition~3.4.9, pp.~181--182; Theorem~3.4.4, p.~180]
{HoveyShipleySmith2000}. Here $i$ runs over boundary inclusions,
$c_j:F_{j+1}S^1\to C\lambda_j$ is the mapping-cylinder
cofibration of $\lambda_j:F_{j+1}S^1\to F_jS^0$, and
$FI_\Lambda$ consists of the free horn inclusions. Their domains
and codomains are finite free-cell spectra. Smashing with $R$
gives finite-cell module maps. The transferred module model
structure is \cite[Theorem~4.1(1), pp.~498--499; proof p.~504]
{SchwedeShipley2000}, applied to $R^{\rm op}$ for right modules.
We also retain the finite cylinder inclusions
$P\to P\wedge\Delta^1_+$.

Take finite pushouts and composites of these acyclic attachments,
and formally invert them. This supplies a finite-certificate description
of every stable equivalence between finite cell modules. To check the
last assertion, construct a fibrant replacement by the small-object
argument using the stated generators, in countably many stages. Any
map from a finite cell module, or from its tensor with a finite
simplicial parameter set, into that replacement factors through a
finite relative cell extension: it uses finitely many generators,
and each attaching map uses finitely many earlier generators.
Induction through the finitely many stages involved gives a finite
set closed under those dependencies. If a map between finite cell
modules is a stable equivalence, its inverse in the model homotopy
category and the two inverse homotopies have precisely such finite
factorizations. They therefore give an inverse after formally
inverting finite acyclic extensions. Conversely every map inverted
by this procedure is a stable equivalence, because all the inverted
attachments are such and cylinder homotopies preserve them. Thus
no geometric concordance test, and no unspecified extra invariant,
is used to select these maps.

Use the hammock localization of this countable category and its
stable equivalences. Its simplices are finite diagrams, so it remains
countable. It computes the derived module mapping spaces:
\cite[Proposition~4.4 and Application~4.6, p.~432]
{DwyerKan1980Functions} gives the function-complex comparison, and
the preceding finite-factorization argument restricts that comparison
to the finite cell objects. It applies also to any finite simplicial
parameter diagram, so verifies both surjectivity and injectivity on
all homotopy groups of the mapping spaces. Select the invertible
components of the endomorphism monoid of $R^{\vee n}$ to obtain
$H_n$. Invertibility itself is witnessed by a finite inverse diagram
and two finite paths to the identity. Direct sum preserves all the
specified attachments and gives the direct-sum coherences on these
models. They can be kept as coherent operations, or rectified by the
simplicial bar construction; that construction uses finite strings
and preserves countability. This specifies the models in (A.119.1)
from finite $K$. It does not assert a terminating procedure for
recognizing a non-equivalence or for solving the resulting group
word problem.

\emph{Acyclic group completion and the two covers.}
Under direct sum the spaces
\[
 \mathcal M=\coprod_{n\geq0}BH_n,
 \qquad \mathcal N=\coprod_{n\geq0}BW_n
\]
are homotopy commutative monoids. The homotopy commutativity is
given by block-permutation conjugation. It is this assertion about
$\mathcal M,\mathcal N$, not about the uncompleted connected space
$C$, that is used here. Their component monoids are $\mathbb N$, and
successively adding one identity summand is cofinal. The
permutation-matrix models in \cite[\S2, p.~247]{Jahren2010}
identify their group completions as
\[
 \Omega B\mathcal M\simeq A(|K|),
 \qquad \Omega B\mathcal N\simeq Q(|K|_+),
                                                               \tag{A.119.4}
\]
and identify their induced map with the unit
$\iota:Q(|K|_+)\to A(|K|)$. Only the indicated rank-zero
components matter; no identification of all projective $K_0$
classes with free rank is being made.

By \cite[Proposition~1.1 and proof pp.~1--7, especially p.~3]
{Gritschacher2026Acyclicity}, the telescope-to-group-completion maps
are acyclic. Its hypotheses hold since the component monoid is
cofinal and central in the homology ring of a homotopy commutative
monoid. Thus the maps from $C$ and $D$ to the corresponding
rank-zero components in (A.119.4) induce abelianization on fundamental
groups and homology isomorphisms with every local coefficient
system. Pull back the simply connected covers of those components.
The pullbacks are exactly $C^{\rm ab}$ and $D^{\rm ab}$, and the
lifted maps are integral homology isomorphisms. Hurewicz on the
simply connected targets gives natural isomorphisms
\[
 H_2(C^{\rm ab};\Z)\cong\pi_2A(|K|),
 \qquad H_2(D^{\rm ab};\Z)\cong\pi_2Q(|K|_+).
                                                               \tag{A.119.5}
\]
The square involving $f^{\rm ab}$ and the lifted unit commutes up
to the specified based homotopy, so these isomorphisms respect the
homomorphism whose cokernel appears in (A.119.3).

\emph{The whole component group, including the extension law.}
The split unit and the stable parametrized h-cobordism theorem give
\[
 A(|K|)\simeq Q(|K|_+)\times\operatorname{Wh}^{\rm Diff}(|K|),
 \qquad
 \pi_0P_{\rm Diff}^{\rm st}(B)
       \cong\pi_2\operatorname{Wh}^{\rm Diff}(|K|).
\]
The split sequence and its unit are read in
\cite[p.~247 and the basic diagram p.~248]{Jahren2010}; the full
geometric comparison is \cite[Theorem~0.1, p.~2; Definition~1.3.4
and (1.3.5), p.~18]{WaldhausenJahrenRognes2013}.
Here the compatibility with the original composition operation
requires a delooping argument. For any simplicial group $P$, the
map $\pi_0P\to\pi_1BP$ takes $g$ to its bar edge; the bar
$2$-simplex $(g,h)$ identifies the edge of the product with the
concatenated edges, using the matching bar convention. Thus this
map is an isomorphism of groups. By
\cite[Remark~1.1.11 and (1.1.12), p.~12]
{WaldhausenJahrenRognes2013}, the component $H(M)_0$ of the
trivial h-cobordism is $BP_{\rm Diff}(M)$. More precisely, the
side-fixed concordance group includes by a group homomorphism and
weak equivalence into the automorphism group of $M\times I$
relative to $M\times0$, whose bar construction classifies that
component. Stabilization on h-cobordisms is the functor
$W\mapsto W\times J$, taking each automorphism $g$ to
$g\times\mathrm{id}_J$ and preserving composition. Hence these
group identifications commute with stabilization. The stabilization
maps of $H$ are simplicial cofibrations
\cite[Definition~1.1.3, p.~8]{WaldhausenJahrenRognes2013}, so
their fundamental groups commute with the colimit. Applying
$\pi_1$ to the based stable h-cobordism equivalence gives
\[
 \operatorname*{colim}_j\pi_0P_{\rm Diff}(M\times J^j)
 \cong\pi_1H^{\rm st}(M)_0
 \cong\pi_1\Omega\operatorname{Wh}^{\rm Diff}(M)
 =\pi_2\operatorname{Wh}^{\rm Diff}(M)
\]
as groups. A based equivalence induces a homomorphism on
$\pi_1$; no preservation of an unlooped multiplication is being
assumed. The split-unit maps and Hurewicz maps are also
homomorphisms. It follows from (A.119.5) that the latter group,
with this precise operation, is exactly (A.119.3), where Hurewicz
carries it to addition of cycles.
The definition (A.119.2) therefore includes all relations between
lifts of quotient classes. It leaves no separate splitting or
extension class to be chosen. Invariance under a change of the
simplicial or spectral models follows from these natural
identifications, rather than from a claim that their sets of
simplices are identical.

\emph{The stated unstable and definable range.}
For a smoothable $n$-dimensional PL manifold choose a smoothing.
Burghelea--Lashof's component isomorphism
\cite[Theorem~6.1, p.~41; proof pp.~42--43]
{BurgheleaLashof1974II}, already checked in A.103, identifies smooth
and PL concordance components. By
\cite[Theorem~1.4.1, p.~21]{WaldhausenJahrenRognes2013}, smooth
stabilization is $k$-connected when
$n\geq\max\{2k+7,3k+4\}$. With $k=1$, every stabilization from
$n\geq9$ onward is $1$-connected, hence bijective on components.
Consequently the unstable group in this range equals the stable
group just presented. A.113 identifies its components with the
definable group over every real closed field. The finite product
statement follows because a concordance fixed on the bottom cannot
permute connected components.

Smoothability is essential to this particular range deduction:
the opened book explicitly separates the general PL stability
problem from its smoothable Corollary~1.4.2 (p.~22). No unstable
conclusion outside the stated class is obtained by this argument.
\end{proof}


\begin{saresult}{A.120}{Proved under stated hypotheses}{Stable components without a smoothing hypothesis}
Let $M$ be any compact PL manifold, with the side boundary fixed in
the definition of concordance, and let $K$ be a finite simplicial
homotopy model for $M$. The complete countable presentation $E(K)$
of A.119 computes its stable PL concordance component group:
\[
 \operatorname*{colim}_{j}\pi_0P_{\rm PL}(M\times I^j)
       \cong E(K).
\]
There is no smoothability or dimension hypothesis on this stable
statement. It does not identify the unstable group with the stable
group for a nonsmoothable PL manifold.
\end{saresult}
\begin{proof}
In this proof Whitehead symbols denote the connective spectra,
rather than only their zeroth infinite-loop spaces. The definitions
in \cite[Definition~1.3.2 and (1.3.3), pp.~16--17;
Definition~1.3.4 and (1.3.5), p.~18]
{WaldhausenJahrenRognes2013} give
\[
\begin{split}
 \operatorname{Wh}^{\rm Diff}(K)
   &=\operatorname{cofib}(\Sigma^\infty K_+\longrightarrow A(K)),\\
 \operatorname{Wh}^{\rm PL}(K)
   &=\operatorname{cofib}(A(*)\wedge K_+\longrightarrow A(K)).
\end{split}
\]
The first arrow factors through the second by the unit
$\mathbb S\to A(*)$. Replace this composable pair by cofibrations
$U\hookrightarrow V\hookrightarrow W$ without changing its
homotopy cofibers. The quotient sequence
$V/U\to W/U\to W/V$ is a cofiber sequence. Therefore the
precise comparison sequence is
\[
 \operatorname{Wh}^{\rm Diff}(*)\wedge K_+
 \longrightarrow\operatorname{Wh}^{\rm Diff}(K)
 \longrightarrow\operatorname{Wh}^{\rm PL}(K).
                                                        \tag{A.120.1}
\]
This derives the comparison from the defining assembly maps and
does not presume that $M$ has a smoothing.

The spectrum $\operatorname{Wh}^{\rm Diff}(*)$ is $2$-connected.
It is connective, the unit is an isomorphism on rank $\pi_0$, and
$\pi_1\operatorname{Wh}^{\rm Diff}(*)=
\operatorname{Wh}_1(1)=0$; see
\cite[Theorem~0.3 and the following paragraph, p.~4]
{WaldhausenJahrenRognes2013}. To check degree $2$ without assuming
an extra algebraic $K$-group calculation, take a smooth disk $D^m$
with $m\geq9$. The group comparison of A.119 identifies
\[
 \pi_2\operatorname{Wh}^{\rm Diff}(*)
 \cong\operatorname*{colim}_{j}
           \pi_0P_{\rm Diff}(D^m\times I^j).
\]
After rounding corners each base is a disk. Burghelea--Lashof's
component comparison, already checked in A.103, identifies each
term with its PL term. Those PL disk concordance component groups
are zero by the Alexander trick, as explicitly used in
\cite[proof of Corollary~1.4.2, p.~23]
{WaldhausenJahrenRognes2013}. Thus the displayed group is zero.

The finite cell filtration of $K_+$ builds the smash product on the
left of (A.120.1) from nonnegative suspensions of a $2$-connected
spectrum. Its homotopy groups in degrees at most $2$ vanish.
The long exact sequence of (A.120.1) consequently gives a natural
group isomorphism
\[
 \pi_2\operatorname{Wh}^{\rm Diff}(K)
       \xrightarrow{\sim}\pi_2\operatorname{Wh}^{\rm PL}(K).
\]
The left side is $E(K)$ by A.119. On the right, the PL version of
the stable parametrized h-cobordism theorem and the bar-space
group-law argument in A.119 identify it with the stable PL
concordance component group. Both comparisons are isomorphisms of
groups. No unstable PL stabilization theorem has been used.
\end{proof}

\paragraph{Coefficient fields and the remaining unstable question.}
The same stable PL group is obtained from a realization over any
real closed field. One can use the relative version of A.115 at
each finite stage. Likewise the component argument in A.113
applies with the entire subpolyhedron
$M\times\{0\}\cup\partial M\times I$ fixed, so the stagewise PL
and definable component groups coincide. The standard PL positive
suspension can be chosen with fixed rational polyhedral bends;
these comparisons commute with it. This also identifies the
stable definable component group with $E(K)$. For a particular
unstable nonsmoothable manifold, the remaining precise issue is
whether its map to this stable group is bijective; A.120 supplies
no such assertion.


\begin{saresult}{A.121}{Proved under stated hypotheses}{A complete presentation for arbitrary PL manifolds in dimension at least nine}
Let $M$ be a compact connected PL manifold of dimension $n\geq9$,
possibly with boundary. All concordances fix the side boundary.
Choose a finite combinatorial triangulation $K$ of $M$, and let
$K^{(3)}$ be its $3$-skeleton. Then the countable abelian group
presentation $E(K^{(3)})$ of A.119, including every cycle,
boundary, and unit-image relation specified there, presents
\[
             \pi_0 P_{\rm PL}(M).
\]
No smoothing of $M$ is assumed. The same presentation gives the
components of the stipulated definable concordance group of any
compact definable manifold with this PL model, over any real closed
field and in any fixed o-minimal expansion. Thus this is an unstable
component computation, whereas A.120 computes the stable group in
every dimension. The presentation is countable; no finite invariant
factor decomposition or solution to its word problem is asserted.
For a disconnected compact manifold, use the finite product of the
presentations of its connected components.
\end{saresult}
\begin{proof}
We first work over $\mathbb R$. If $M$ has boundary, choose a PL collar
and push its outside face a positive distance into the collar. This
gives a PL embedding $e:M\longrightarrow\operatorname{int}M$
homotopic, as a map to $M$, to the identity. For example, on a collar
$\partial M\times[0,2]$ replace its second coordinate $t$ by
$\epsilon+(1-\epsilon/2)t$, with $0<\epsilon<1$, and leave the
complement of the collar fixed. The formulas agree at $t=2$;
replacing $\epsilon$ by $s\epsilon$ supplies the homotopy.
If $M$ is closed, take $e=\operatorname{id}$.
Put $L=e(|K^{(3)}|)$ and choose a sufficiently small closed regular
neighborhood $N$ of $L$ in $\operatorname{int}M$.

Here are the required properties of this neighborhood. They can all
be read from a finite subdivision, so no limiting neighborhood is
being substituted for $N$. Triangulate the pair $(M,L)$ with a mesh
small enough near $L$, and take its first derived subdivision, in
which the subdivision of $L$ is full.
If $\lambda$ is the sum of barycentric coordinates at its vertices,
the regular neighborhood may be taken to be $\lambda\geq1/2$ after
subdividing the cut cells. Shrinking inside the chosen open
neighborhood assures $N\subset\operatorname{int}M$. Its boundary
is $\lambda=1/2$, and it is a compact PL $n$-manifold. These are
the derived-neighborhood and simplicial-neighborhood constructions
of \cite[Chapter~3, Proposition~3.10 and Theorem~3.11, pp.~34--35,
and Corollary~3.18, p.~36]
{RourkeSanderson1972}.

Within each simplex write a point with $0<\lambda<1$ as
$\lambda u+(1-\lambda)v$, where $u$ uses its $L$-vertices and
$v$ its other vertices. Moving $\lambda$ to $1$ retracts $N$ onto
$L$. Moving it to $1/2$ retracts $N\setminus L$ onto $\partial N$.
The barycentric formulas agree on common faces, and are continuous
on their respective domains. In particular $N$ has the homotopy
type of the finite $3$-complex $K^{(3)}$ and is connected.

The inclusion $K^{(3)}\longrightarrow K$ is $3$-connected, since
all the omitted cells have dimension at least four. The retraction
and the homotopy from $e$ to the identity therefore show that
\[
                 \pi_i(M,N)=0\qquad (i\leq3).
                                                        \tag{A.121.1}
\]
Moreover, relative PL general position moves maps of intervals,
circles, and disks of dimension at most two off $L$, keeping any
part already in $N\setminus L$ fixed. The dimension inequality is
$2+3<n$. One may first push the portion of a map not fixed at
$\partial N$ into $\operatorname{int}N$ by a collar. The usual
relative general-position argument then applies there, and the
collar does not meet $L$; see \cite[Chapter~5, Theorems~5.3--5.4]
{RourkeSanderson1972}. For clarity in the injectivity argument,
represent a boundary loop by an embedded PL loop in $\partial N$,
which is possible by general position, and choose an embedded collar
on that loop into $N\setminus L$. Apply Theorem~5.4 to a filling
disk relative to this collar; its dimension bound makes the disk an
embedding. Then Theorem~5.3 moves its image off $L$ while fixing the
collar. Thus the relative data in the map theorem are indeed PL and
nondegenerate. The interval and circle arguments are the same in
lower dimension. Consequently
$N\setminus L\longrightarrow N$ induces isomorphisms on
$\pi_0$ and $\pi_1$. With the preceding boundary retraction this
proves
\[
       \pi_i(\partial N)\xrightarrow{\ \cong\ }\pi_i(N)
                      \qquad (i=0,1).                  \tag{A.121.2}
\]

The same regular-neighborhood construction gives $N$ a handle
decomposition with handles of indices at most three. Namely, thicken
the simplices of $L$ successively in increasing dimension. The
normal block of an $r$-simplex in the ambient PL manifold is a disk,
so the thickening step attaches $D^r\times D^{n-r}$ along
$\partial D^r\times D^{n-r}$. The product structures agree on
the previously thickened faces. This is the finite handle
construction from a triangulation; compare
\cite[Proposition~6.9 and its proof, pp.~82--83]
{RourkeSanderson1972}.
It also proves directly that $N$ has a compatible smooth structure.
Start with its $0$-handles and apply the classical PL--DIFF
handle-smoothing argument. The obstruction to extending over a
$k$-handle is the concordance class of the induced smoothing of
$S^{k-1}$; for $k=1,2,3$ this class is zero. The local product
structure and concordance-implies-isotopy theorems make the extension
agree with the already chosen smoothing near the attaching region.
This precise handle argument is given in
\cite[Essay~IV, Appendix~B, p.~211]{KirbySiebenmann1977}; it supplies
a smoothing compatible with the given PL structure, rather than
merely a smoothing of the underlying topological space. In
obstruction-theoretic language, the same conclusion follows by
lifting the PL tangent classifying map over a $3$-dimensional spine;
only the low, vanishing obstruction groups occur.

We now apply the PL localization theorem of Burghelea, Lashof and
Rothenberg. Its exact form is restated, with its
extension-by-identity map, in
\cite[Theorem~3.3, p.~17]{GalvinNonino2025}, citing
\cite[Theorem~3.1$'$]{BurgheleaLashofRothenberg1975}.
For $W\subset\operatorname{int}V$, it requires
$\pi_i(\partial W)\cong\pi_i(W)$ for $i=0,1$,
$\pi_i(V,W)=0$ for $i\leq r$, and $W$ to be $k$-connected,
where $0\leq k\leq r\leq\dim V-4$. It says that the map
\[
       \gamma:P_{\rm PL}(W)\longrightarrow P_{\rm PL}(V),
       \qquad \gamma(h)=h\text{ on }W\times I,
\]
extended by the identity elsewhere, has zero relative homotopy
groups in degrees at most
$\min(2r-3,r+k-2,r+2)$.
The relative pair is $(V,W)$: this agrees with the original
Theorem~3.1 on printed p.~28, which has been opened and visually
checked; the displayed $(W,V)$ in the later restatement reverses
the pair. The original pp.~28--29 establish the precursor theorem
and its handle proof. The later Theorem~3.1$'$ is used through the
explicit, opened restatement just cited; its original page has not
been retrieved.

Take $V=M$, $W=N$, $r=3$, and $k=0$. Equations
(A.121.1)--(A.121.2) check the hypotheses, and $3\leq n-4$.
The bound is
\[
                   \min(3,1,5)=1.
\]
Thus $\gamma$ induces a bijection on path components.
It preserves composition on the nose, because both extensions are
the identity off $N\times I$. Its component bijection is therefore
a group isomorphism
\[
          \pi_0P_{\rm PL}(N)\xrightarrow{\ \cong\ }
          \pi_0P_{\rm PL}(M).                          \tag{A.121.3}
\]
This step uses the actual localization map, rather than an unmarked
homotopy equivalence of the underlying spaces.

The manifold $N$ is smoothable, has dimension $n\geq9$, and has
finite homotopy model $K^{(3)}$. Apply A.119, including its boundary
version, to obtain
$\pi_0P_{\rm PL}(N)\cong E(K^{(3)})$ as groups.
Together with (A.121.3) this proves the asserted full presentation.
It is unnecessary here to identify this presentation with
$E(K)$ or to assert a PL stabilization theorem for $M$.

Finally take a compact definable manifold over a real closed field
and a finite PL model $K$. Realize the same abstract triangulation
over the real algebraic numbers and hence over $\mathbb R$.
A.115 identifies the PL concordance components over these fields,
and the relative, component comparison A.113 identifies them with
the definable components in any fixed o-minimal expansion. Both
comparisons preserve composition. Thus the real group presentation
above gives exactly the stipulated definable component group as
well. No real model of the o-minimal expansion is used. A concordance
fixing the bottom cannot permute the finitely many connected
components of $M$, which proves the final product assertion.
\end{proof}


\begin{saresult}{A.122}{Proved under stated hypotheses}{The complete component presentation in every PL dimension at least six}
Let $M$ be a compact connected PL manifold of dimension $n\geq6$,
with any boundary fixed on the sides of a concordance, and let $K$
be a finite simplicial model of its homotopy type. Then
\[
                   \pi_0P_{\rm PL}(M)\cong E(K),
\]
where $E(K)$ is the complete countable group presentation of
A.119. The same formula holds for the stipulated definable
concordance group of a compact definable manifold with this PL
model, over every real closed field and in every fixed o-minimal
expansion. Use the product of these presentations for finitely many
connected components, each of dimension at least six.

More precisely, positive suspension induces an isomorphism on
components of smooth concordance groups once $n\geq6$. This
identifies the unstable group of a smoothable PL model with the
stable presentation. The nonsmoothable PL case is obtained by
localizing to a smoothable neighborhood of a $3$-skeleton, not by
assuming a general unstable PL stability theorem. No assertion
about arbitrary manifolds of dimensions four or five is included.
\end{saresult}
\begin{proof}
We first lower the dimension bound for the smoothable case, keeping
track of the full relation subgroup in the corrected obstruction
sequence. For a connected space $X$ write
\[
 A_X=\operatorname{Wh}_1^+(\pi_1X;\mathbb Z/2\oplus\pi_2X),
 \qquad
 I_X=\operatorname{im}\bigl(\chi_X:
       \operatorname{Wh}_3(\pi_1X)\longrightarrow A_X\bigr).
\]
Replacing $\operatorname{Wh}_3$ by its surjecting algebraic
$K_3(\mathbb Z[\pi_1X])$ does not change $I_X$.
We use the corrected geometric exact-sequence theorem, in its full
form, as stated in
\cite[Theorem~3.1, p.~6]{Igusa2021Dimension3}: for a compact smooth
manifold of dimension at least six, it gives
\[
 0\longrightarrow A_M/I_M\xrightarrow{\ \overline\theta_M\ }
 \pi_0P_{\rm Diff}(M)\xrightarrow{\ \Sigma_M\ }
 \operatorname{Wh}_2(\pi_1M)\longrightarrow0.             \tag{A.122.1}
\]
The $\chi$ here is the full corrected obstruction map, including
its framing component. It is the same map as in the stable
corrected sequence in \cite[Introduction, (0.2), p.~2]
{Igusa2021SecondObstructionI}, and in the functorial A-theory
construction of \cite[Section~2, especially Corollary~2.7, p.~250]
{Jahren2010}. This identification of the full corrected sequence is
a cited input. In particular, we are not deriving it from the
formula for the $\pi_2M$-component of $\chi$ printed after Igusa's
Theorem~3.1. That partial formula would not determine $I_M$.
The original 1984 chapter \cite{Igusa1984Postnikov} has not been
retrieved; the precise opened restatements and functorial
construction just listed are the sources used here.

The deformation retraction $M\times I\longrightarrow M$
identifies $\pi_1$, its action on $\pi_2$, and their first Postnikov
invariant. Functoriality of the full sequence in
\cite[Corollary~2.7]{Jahren2010} identifies $A_{M\times I}$ with
$A_M$ and $I_{M\times I}$ with $I_M$. We next verify that these
are the vertical identifications induced by positive suspension on
the two geometric obstruction maps in (A.122.1).

Use concordances that are the identity near the bottom and sides,
and are products near the top. Inclusion of this normalized
subgroup induces the same component group; collar straightening
also works for compact one-parameter families. The positive
suspension formula of
\cite[Section~4.1, p.~8]{Igusa2021Dimension3} applies a
concordance $g$ on each fixed polar ray and is the identity on the
remaining region. Thus
\[
                     \sigma^+(gh)=\sigma^+(g)\sigma^+(h)
\]
on that subgroup. It consequently gives a homomorphism of
component groups, respecting the bottom and side collars.

For the second obstruction use marked lens models in the fixed
indices $3,4$ in every dimension $n\geq6$. This convention avoids
an index change when $n$ changes parity. The two-index deformation
is available in this range: the original
\cite[Part~I, Chapter~V, Theorem~3.1(a), p.~184]
{HatcherWagoner1973} permits $2\leq i\leq n-2$, and its use for
the one-parameter representatives is explicit in Part~II, Section~2,
p.~244. The marked-lens framed-bordism calculation of
\cite[Theorem~1.8 and Lemma~1.10, pp.~8--9]
{Igusa2021SecondObstructionI} applies with $k=3$: its proof needs
the inequalities $k\geq3$ and $n-k\geq3$, which hold here.
It identifies the marked-lens group with
\[
 (\mathbb Z/2\oplus\pi_2M)[\pi_1M],/
       (\mathbb Z/2\oplus\pi_2M)[1].
\]
Changing the markings gives the coinvariants $A_M$; forgetting the
lens model gives $\theta_M$. Use the stable sign convention
$\lambda_0^s=(-1)^k\lambda_0$ from
\cite[Section~4.2, pp.~8--9]{Igusa2021Dimension3}.
That source explicitly proves compatibility with positive
suspension, which retains $k$, for the complete framed invariant.
This includes its $\mathbb Z/2$ coordinate. Hence the square
\[
 \sigma^+_*\theta_M=\theta_{M\times I}
       : A_M\longrightarrow\pi_0P_{\rm Diff}(M\times I)
                                                               \tag{A.122.2}
\]
commutes under the preceding identification of $A$.
The first obstruction is unchanged under the same suspension:
\cite[Part~I, Chapter~IV, Section~1, Remark, p.~130]
{HatcherWagoner1973}, using the trajectory comparison in
Chapter~I, Section~5, (5.2), p.~60. Therefore
\[
       \Sigma_{M\times I}\sigma^+_* =\Sigma_M.          \tag{A.122.3}
\]

Equations (A.122.2)--(A.122.3) give a commutative diagram of
short exact sequences (A.122.1), whose left and right vertical
arrows are the identity on $A_M/I_M$ and
$\operatorname{Wh}_2(\pi_1M)$, respectively. The middle arrow is
$\sigma^+_*$. It is an isomorphism by exactness. Explicitly, if
$\sigma^+_*(g)=1$, then $\Sigma_M(g)=0$, so
$g=\overline\theta_M(a)$; injectivity of the lower
$\overline\theta$ forces $a=0$ and hence $g=1$.
Given $g'$ in the target, first choose $h$ with
$\Sigma_M(h)=\Sigma_{M\times I}(g')$. The difference
$g'\sigma^+_*(h)^{-1}$ lies in the image of
$\overline\theta_{M\times I}$ and is therefore the suspension of
an element of $\operatorname{im}\overline\theta_M$. This proves
surjectivity. This argument does not split either extension.

The compact-manifold and boundary hypotheses persist throughout.
In \cite[Introduction, p.~1]{Igusa2021SecondObstructionI},
$P_{\rm Diff}(M)$ is explicitly
$\operatorname{Diff}(M\times I\text{ rel }M\times0\cup
\partial M\times I)$. Section~1, p.~3, fixes the projection
function near the entire boundary. The lens modifications and
two-index deformations are supported in the interior. The
Hatcher--Wagoner setup likewise uses compact pairs $(M,\partial M)$,
and the collar normalization preceding suspension is on printed
p.~55. Thus the argument applies to every $M\times I^j$ after
rounding its corners. Iterating the component isomorphism
identifies $\pi_0P_{\rm Diff}(M)$ with its stable group once
$n\geq6$. The complete stable group presentation in A.119 is
$E(K)$, with its group law already checked through classifying
h-cobordisms. The smooth-to-PL component isomorphism of A.103
therefore proves the desired statement whenever $M$ is smoothable
and $n\geq6$.

Now let $M$ be any compact PL manifold of dimension $n\geq7$.
Repeat the finite regular-neighborhood construction of A.121.
All of its localization hypotheses require only $3\leq n-4$ and
$2+3<n$, which hold here. The resulting $N$ is a smoothable
$n$-manifold with the homotopy type of the chosen $3$-skeleton.
The just-proved smoothable result now applies to $N$ in this lower
dimension. Extension by the identity consequently gives
\[
                 \pi_0P_{\rm PL}(M)\cong E(K^{(3)}).
                                                               \tag{A.122.4}
\]
Here first take $K$ to be a triangulation of $M$.
There is a natural group isomorphism
$E(K^{(3)})\longrightarrow E(K)$. To see it, use the functorial
sequence of \cite[Corollary~2.7]{Jahren2010} for the inclusion
$K^{(3)}\longrightarrow K$. This inclusion is an isomorphism on
$\pi_1$ and $\pi_2$, so it induces isomorphisms on
$\operatorname{Wh}_3$, $A$, and $\operatorname{Wh}_2$.
The square with the full $\chi$ commutes, so it also induces an
isomorphism on its cokernel. Applying exactness to
\[
 0\longrightarrow A_X/I_X\longrightarrow
 \pi_2\operatorname{Wh}^{\rm Diff}(X)\longrightarrow
 \operatorname{Wh}_2(\pi_1X)\longrightarrow0
\]
therefore gives an isomorphism on the middle groups. A.119 identifies
these groups with $E(K^{(3)})$ and $E(K)$, proving the claim.
This also allows any finite simplicial homotopy model $K$ in the
statement, rather than only the selected triangulation.

For completeness, the remaining dimension six case needs no
localization: every PL $6$-manifold admits a compatible smoothing.
Here is the exact classical input and its application. A finite PL
handle decomposition has handle indices at most six. The
PL--DIFF extension argument in
\cite[Essay~IV, Appendix~B, p.~211]{KirbySiebenmann1977}
identifies the obstruction at a $k$-handle with the class of a
smooth twisted $(k-1)$-sphere, namely an element of
$\Gamma_{k-1}$. These groups vanish for $0\leq k-1\leq5$.
For the low-dimensional compatible-smoothing assertion one can
read \cite[Essay~I, p.~39]{KirbySiebenmann1977}; the remaining
values $\Gamma_4=0$ and $\Theta^{\rm Diff}_5=0$ are explicitly
stated in Essay~IV, Section~10.12, p.~200. A smooth twisted
$5$-sphere represents a smooth homotopy $5$-sphere, so the latter
equality gives $\Gamma_5=0$. Starting with the $0$-handles and
extending across each successive handle now produces a compatible
smooth structure on $M$, including a smooth structure on its
boundary. No preassigned boundary smoothing is required here.
The smoothable case already proved applies and gives $E(K)$.

Finally A.113 and A.115 transfer the PL component calculation to
the stipulated definable concordance group over any real closed
field. Their relative versions preserve the side-boundary
convention and composition. Finite products handle disconnected
manifolds as in A.121. The lower-dimensional exclusions in the
statement remain: neither the corrected exact-sequence input nor
the fixed-index argument used here asserts this conclusion for
arbitrary $4$- or $5$-manifolds.
\end{proof}


\begin{saresult}{A.123}{Proved under stated hypotheses}{Trivial low-dimensional component groups and the ball case}
Let $P_{\rm PL}(B)$ fix $B\times\{0\}\cup\partial B\times I$
pointwise, with the upper face free. Then
\[
 \pi_0P_{\rm PL}(B)=0
\]
if $B$ is a compact manifold of dimension zero or one, if $B$ is
a PL ball of any dimension, if $B$ is a compact orientable
surface other than $S^2$, or if $B$ is a closed nonorientable
surface other than $\mathbb{RP}^2$. The same assertion holds for a finite
disjoint union of these manifolds and for their definable models
over every real closed field, in the stipulated component convention.
No higher homotopy-group assertion is made here. The two closed
surfaces $S^2$ and $\mathbb{RP}^2$, and nonorientable surfaces
with boundary, are not included in the surface clauses.
\end{saresult}
\begin{proof}
All isotopies in this proof are jointly PL. In particular,
replacing an isotopy parameter $u$ by a product $st$ is unnecessary.
We use instead the PL function
\[
 \rho(s,t)=\max\{0,s+t-1\},\qquad (s,t)\in I^2.
                                                               \tag{A.123.1}
\]
It is zero if $s=0$ or $t=0$, equals $t$ if $s=1$, and equals
$s$ if $t=1$. Composition with a finite PL family therefore
remains a finite PL map, while preserving every parameter that
is explicitly retained as an output coordinate.

First suppose $B=D^m$. A concordance $h$ restricts on the upper
face to a PL homeomorphism $a$ of $D^m$ fixed on its boundary.
The PL Alexander isotopy gives a jointly PL path $a_u$ from the
identity to $a^{-1}$, fixed on $\partial D^m$; the explicit
level-preserving cone construction is given in
\cite[Theorem~5.2, pp.~95--96]{Epstein1966}.
For completeness, for $m\geq1$ on the cube $D^m=[-1,1]^m$ it is
\[
 a_u(x)=
 \begin{cases}
 u\,a^{-1}(x/u),&\|x\|_\infty\leq u,\ u>0,\\
 x,&\|x\|_\infty\geq u,
 \end{cases}
 \qquad a_0(x)=x.
\]
The two formulas agree on their common boundary. If
$a^{-1}(y)=Ay+b$ on a cell, the first formula is $Ax+ub$
on the cone over that cell at height one. These are finitely
many polyhedral cells, so the total isotopy is PL in every
dimension; the same formula using $a$ gives its inverse.
The maps
\[
 A_s(x,t)=(a_{\rho(s,t)}(x),t)
\]
are concordances, start at the identity, and make $A_1h$ the
identity on the entire boundary of $D^m\times I$.
This cylinder is a PL ball, so a second relative Alexander
isotopy deforms $A_1h$ to the identity. Concatenating the two
isotopies proves the ball assertion, including $m=0$ and $m=1$.
The displayed map and its inverse are continuous and PL: each
slice is a homeomorphism, and the total map is a continuous
bijection of compact polyhedra. Equivalently the inverse is
obtained by substituting the inverse PL isotopy family.

It remains in dimension one to treat $B=S^1=\mathbb R/\mathbb Z$.
Write $A=S^1\times I$. A concordance is orientation preserving
and its upper-circle restriction has degree one, since it is
homotopic to the bottom restriction through the cylinder.
An orientation-preserving PL circle homeomorphism is joined to
the identity by a PL isotopy: lift it to an increasing PL map
$\widetilde a:\mathbb R\to\mathbb R$ satisfying
$\widetilde a(x+1)=\widetilde a(x)+1$, triangulate the finitely
many breakpoint rectangles in $[0,1]\times I$, and linearly
interpolate their bottom and top vertex values. The horizontal
slopes on the two triangles of each rectangle are respectively
the positive slopes of $\widetilde a$ and of the identity.
Thus every horizontal slice is increasing, and the periodic
boundary values agree. This gives the required circle isotopy.
Using (A.123.1) extends an inverse upper-circle isotopy across
$A$, fixed on its bottom. We may consequently assume $h$ fixes
both boundary circles.

Let $\alpha=\{0\}\times I$. The proper arc $h(\alpha)$ has
an integral winding number $k$ relative to its fixed endpoints,
computed in the strip cover $\mathbb R\times I$.
The PL homeomorphism
\[
 T_{-k}([x],t)=([x-kt],t)
\]
fixes both boundary circles and makes $T_{-k}h(\alpha)$
homotopic to $\alpha$ relative to the endpoints. Moreover
$T_{-k}$ lies in the identity component of the concordance
group: its explicit isotopy is
\[
 T_{-k,s}([x],t)=([x-k\rho(s,t)],t).
\]
Only the bottom circle is required to stay fixed during this
isotopy. Epstein's relative arc theorem
\cite[Theorem~3.1, p.~89]{Epstein1966} supplies a PL ambient
isotopy fixed on both boundary circles taking the adjusted arc
to $\alpha$. Its interval parametrization can also be fixed,
using the relative interval Alexander isotopy in a regular
neighborhood, exactly as in the proof of
\cite[Theorem~3.4, p.~91]{Epstein1966}.
Cutting $A$ along $\alpha$ now gives a PL disk homeomorphism
fixed on its entire boundary. The relative Alexander isotopy
on this disk descends to the annulus. This proves the circle
assertion without retaining a Dehn-twist obstruction: the
upper face was free during the displayed $T_{-k,s}$ isotopy.

Finally let $F$ be a compact connected orientable surface
different from $S^2$. Waldhausen explicitly works in the PL
category, and in Section~3 assumes precisely that $F$ is
orientable and is not the sphere
\cite[pp.~56,65]{Waldhausen1968}. His Lemma~3.5, p.~69,
states that a homeomorphism of $F\times I$ fixed on
$F\times\{0\}\cup\partial F\times I$ is isotopic, keeping
the entire boundary fixed, to a level-preserving homeomorphism.
Apply the lemma to a concordance $h$, and write the resulting
homeomorphism as $(x,t)\mapsto(h_t(x),t)$, where $h_0$ is the
identity and every $h_t$ fixes $\partial F$.
Then
\[
 (x,t,s)\longmapsto(h_{\rho(s,t)}(x),t,s)
\]
is a jointly PL isotopy through concordances from the identity
to this level-preserving map. Its inverse is obtained from the
inverse level-preserving family. Composing with Waldhausen's
isotopy proves the surface clause, including the torus and all
orientable surfaces with boundary.

For completeness, closed nonorientable surfaces other than
$\mathbb{RP}^2$ also satisfy the assertion; this step uses the
nonorientable three-dimensional isotopy theorem, rather than
extending Waldhausen's orientable hypothesis without proof.
Let $F$ be such a surface and $V=F\times I$.
The universal cover of $F$ is the plane and its fundamental
group is torsion free (see the surface facts in
\cite[Corollary~1.8, pp.~85--86, and Lemma~4.3, p.~92]{Epstein1966}),
so the universal cover of $V$ is
$\mathbb R^2\times I$. An embedded PL two-sphere in $V$
lifts to a two-sphere bounding a PL ball in this slab, by the
three-dimensional PL Schoenflies theorem. Distinct translates
of its lifted sphere are disjoint. Their bounded balls are
either disjoint or nested. Proper discontinuity excludes
nesting: iteration of the corresponding nontrivial, infinite
order deck transformation would place infinitely many points
of an orbit in one compact ball. Thus the ball projects
injectively, and $V$ is irreducible. A two-sided embedded
$\mathbb{RP}^2$ cannot occur: its fundamental group maps
trivially to the torsion-free $\pi_1V$, so it would lift to
the orientable slab, where a two-sided embedded surface must
be orientable. Finally $F\times\{1/2\}$ is a two-sided
incompressible surface, since its inclusion is a homotopy
equivalence. It cannot bound a contractible submanifold:
otherwise its nontrivial injective fundamental-group map to
$V$ would factor through the zero fundamental group of that
submanifold. Hence $V$ is $\mathbb{RP}^2$-irreducible and
sufficiently large in exactly the definitions of
\cite[Chapter~II, Section~1, pp.~15--17]{Laudenbach1974}.

Write a concordance as $h(x,t)=(u(x,t),v(x,t))$.
It is properly homotopic to the identity. Indeed, first use
\[
 (x,t)\longmapsto
 \bigl(u(x,t),(1-s)v(x,t)+st\bigr),
\]
and then use
\[
 (x,t)\longmapsto\bigl(u(x,(1-s)t),t\bigr),\qquad 0\leq s\leq1.
\]
Here ``properly'' is the boundary-pair convention: throughout,
each of the two boundary components maps into itself, and
interior points remain interior. These formulas are only
homotopies of maps; they need not be isotopies.
Laudenbach's Theorem~II.6.1, pp.~46--47, deforms a proper
homotopy of a homeomorphism to the identity to an isotopy in
this class of manifolds. His category comparison in Chapter~I, Section~1,
pp.~7--9, first gives a TOP isotopy of $h$ to a smooth map $g$.
The map $g$ is still properly homotopic to the identity, so the
smooth formulation of Theorem~II.6.1 gives a TOP isotopy of
$h$ to the identity. The injectivity of
$\pi_0\operatorname{Aut}_{\rm PL}(V)\to
 \pi_0\operatorname{Aut}_{\rm TOP}(V)$ on pp.~7--8 then gives
a jointly PL isotopy $H_s:V\to V$ between the original PL
endpoints. This uses no assertion that the auxiliary smooth
map is literally PL in a fixed triangulation.

It remains to enforce the fixed bottom, since no unmentioned
relative version of that theorem is needed. Every $H_s$
preserves each of the two boundary components, by continuity
from its initial value. Let $b_s$ be its restriction to the
bottom, regarded as a PL homeomorphism of $F$.
The inverse total boundary family is PL, so
\[
 \widehat H_s=(b_s^{-1}\times\operatorname{id}_I)\,H_s
\]
is a jointly PL isotopy through concordances.
Both $b_0$ and $b_1$ are the identity. Thus this isotopy has
the same endpoints $h$ and the identity and fixes the bottom
pointwise at every time. This proves the additional closed
nonorientable clause, including the Klein bottle.

Fixing the bottom prevents permutation of components, so finite
disjoint unions give finite products of the groups just treated.
All real-field proofs consist of finite PL data. A.115 and the
relative component argument of A.113 transfer them to any real
closed field and identify these PL components with the
definable components. For manifolds with boundary the same
relative argument fixes the additional subpolyhedron
$\partial B\times I$ throughout.
\end{proof}


\begin{saresult}{A.124}{Proved under stated hypotheses}{The component presentation in dimension five when the second homotopy group vanishes}
Let $M$ be a compact connected PL manifold of dimension five,
possibly with boundary, and assume $\pi_2(M)=0$.
For every finite connected pointed homotopy model $K$ of $M$,
\[
 \pi_0P_{\rm PL}(M)\cong E(K),
\]
where $E(K)$ is the complete countable group presentation in
A.119. The same calculation holds for the stipulated definable
concordance components of any definable realization over a real
closed field. It determines the full group, including its
extension law; it is not a finite invariant-factor algorithm.
There is no assertion here for five-manifolds with nonzero
$\pi_2$, or for higher homotopy groups of the comparison.
\end{saresult}
\begin{proof}
First give $M$ a compatible smoothing; its existence will be
verified below. Put $\pi=\pi_1M$ and
\[
 A=\operatorname{Wh}_1^+(\pi;\mathbb Z/2)
   =H_0(\pi;(\mathbb Z/2)[\pi])/(\mathbb Z/2)[1],
 \qquad
 I=\operatorname{im}\bigl(
 \chi_{\rm Wh}:\operatorname{Wh}_3(\pi)\longrightarrow A\bigr).
\]
We invoke the full geometric exact-sequence theorem
\cite[Theorem~7.1, p.~23]{Igusa2025}, which explicitly assumes
$\dim M\geq5$ and $\pi_2M=0$ and includes compact manifolds
with boundary, with the bottom and side fixed. It gives
\[
 0\longrightarrow A/I\xrightarrow{\theta_M}
 \pi_0P_{\rm Diff}(M)\xrightarrow{\Sigma_M}
 \operatorname{Wh}_2(\pi)\longrightarrow0.              \tag{A.124.1}
\]
Exactly the same theorem applies to $M\times I$, after rounding
corners. Projection identifies its fundamental group with
$\pi$, and its second homotopy group is still zero. The map
$\chi_{\rm Wh}$ and hence $I$ are the identical group-theoretic
map and subgroup in these two sequences; naturality is proved
in \cite[Section~4, pp.~14--16]{Igusa2025}.
We use Theorem~7.1 as a cited full theorem. The paper explicitly
describes its geometric proof as partial and refers to the
corrected obstruction theory for the remaining details; we do
not infer this full-kernel assertion from a partial coefficient
calculation.
Two typographical errors in that source are kept separate from
its theorem: in Theorem~7.1(a), p.~23, the printed
$H_2\operatorname{GL}(\mathbb Z[\pi])'\cong K_3(\mathbb Z[\pi])$
has $K_3$ in place of $K_2$; the introductory paragraph on
p.~2 says $\pi_1P$ where its displayed sequence and Theorem~7.1
correctly use $\pi_0P$. We use the standard
$\operatorname{Wh}_2$ and the explicitly displayed component
sequence, consistently with A.119 and A.122.

Consider the actual positive suspension
\[
 \sigma^+:P_{\rm Diff}(M)\longrightarrow
                   P_{\rm Diff}(M\times I).
\]
As in A.122, first normalize concordances on fixed boundary
collars. On the bent product, positive suspension applies the
original concordance on each fixed polar ray. It therefore
preserves composition; the collar normalization is an
equivalence on components. This gives a group homomorphism
$\sigma^+_*$, with the bottom and the entire new side boundary
fixed. The construction is described in
\cite[Section~4.1, p.~8]{Igusa2021Dimension3}.

We check both obstruction maps, rather than identifying only
the abstract end groups in (A.124.1).
The first obstruction is unchanged by suspension:
\[
 \Sigma_{M\times I}\sigma^+_* =\Sigma_M.               \tag{A.124.2}
\]
This is stated in \cite[Part~I, p.~130, following
Section~I.5(5.2), p.~60]{HatcherWagoner1973}; adding a positive
quadratic coordinate leaves the ordered handle-slide matrices
and their labels unchanged.

For the second obstruction it suffices to consider
$\ker\Sigma_M$. Use a one-parameter Morse representative in
the fixed adjacent indices $2,3$. The opened
\cite[Part~I, Chapter~V, Theorem~3.1(a), p.~184;
Part~II, Section~2, p.~244]{HatcherWagoner1973}
permits $2\leq i\leq n-2$, so $i=2$ is permitted when $n=5$.
The corrected geometric invariant in
\cite[Section~7, items~(1)--(12), p.~24]{Igusa2025}
is obtained as follows. The labelled intersection graph $J$
is cut and reassembled using the inverse handle-slide matrix;
for each nonidentity label $u\in\pi$ this gives a closed
framed one-manifold $J_u$. Its framed bordism class is an
element of $\Omega_1^{\rm fr}=\mathbb Z/2$. Summing these
classes with their labels and then passing to conjugacy
coinvariants and to $A/I$ is the inverse of $\theta_M$ on
$\ker\Sigma_M$ in (A.124.1).

In the positive quadratic suspension, the intersections,
their paths to the basepoint, their labels, and the
handle-slide matrices are unchanged. Each $J_u$ is the same
one-manifold in the equatorial copy of the old intersection
space; its normal framing has one additional trivial normal
coordinate. The framed bordism class is unchanged by this
stabilization. Any sign caused by a convention on ordering
normal vectors is also invisible in $\mathbb Z/2$.
The prescribed agreement of framings at birth-death and
handle-addition points is preserved by appending the same
constant normal vector. Thus the entire construction, including
the corner smoothing of $J_u$, commutes with suspension.
Consequently
\[
 \sigma^+_*\theta_M=\theta_{M\times I}.                \tag{A.124.3}
\]
This step uses only framed one-manifolds and their stable
framing classes. It does not invoke the $i,n-i\geq3$
realization lemma for an arbitrary $\pi_2M$ coefficient;
that coefficient is zero under the present hypothesis.
The realization and full kernel needed here are exactly those
already supplied by the cited dimension-five Theorem~7.1.

Equations (A.124.2)--(A.124.3) give a commutative diagram of
the two short exact sequences (A.124.1), with identity maps
on $A/I$ and $\operatorname{Wh}_2(\pi)$. The middle map is
therefore an isomorphism of groups. Explicitly, a suspended
zero first has zero first obstruction, lies in the image of
$\theta_M$, and has zero coefficient there, proving
injectivity. For surjectivity, lift the first obstruction of
a target element to a source element; their difference lies
in the target $\theta$ image, which is also a suspended image.
No splitting of (A.124.1) has been used.

The target $M\times I$ is six-dimensional, so A.122 identifies
its component group with the stable group $E(K)$, since it
has the same homotopy model as $M$. The preceding group
isomorphism proves the smooth dimension-five formula and
preserves the extension law in the countable presentation.

For a PL five-manifold choose a finite PL handle decomposition,
including collars when its boundary is nonempty. Extend a
compatible smoothing over handles in increasing index.
The classical PL--DIFF handle argument in
\cite[Essay~IV, Appendix~B, p.~211]{KirbySiebenmann1977}
places the index-$k$ obstruction in $\Gamma_{k-1}$.
All indices are at most five. The low-dimensional compatible
smoothings in \cite[Essay~I, p.~39]{KirbySiebenmann1977}
and the explicit equality $\Gamma_4=0$ in
\cite[Essay~IV, p.~200]{KirbySiebenmann1977}
therefore remove every obstruction. This constructs a
compatible smoothing of the given PL manifold; no smoothing
has been prescribed independently on its boundary.
The smooth-to-PL component comparison in A.103, valid in
dimension at least five and relative to the side boundary,
transfers the smooth calculation. Finally A.113 and A.115,
with the same relative boundary convention, transfer the
component group to every definable realization and every
real closed field. The finite model of $M$ in A.119 may be
replaced by any finite pointed homotopy model by the
functorial group identification proved there.
\end{proof}

% Root and lie_theory independent audits PASS; promoted 2026-11-19.
\begin{saresult}{A.125}{Proved under stated hypotheses}{The two-sphere has trivial concordance component group}
For the standard PL two-sphere and every definable model of it over
any real closed field,
\[
 \pi_0P_{\rm PL}(S^2)=\pi_0P_{\rm def}(S^2)=0.
\]
Here the bottom is fixed pointwise and the upper boundary is free.
This is a component assertion, not a higher parametric comparison.
\end{saresult}
\begin{proof}
We first prove the smooth assertion and then make the category and
bottom corrections explicit. Let $D$ be the closed unit three-ball,
let $D_\epsilon$ be a concentric smaller ball, and let $j$ be its
inclusion. Equip all smooth mapping spaces with their $C^\infty$
topologies, and put
\[
 G=\operatorname{Diff}^+(D),\qquad
 E=\operatorname{Emb}^+(D_\epsilon,\operatorname{int}D),\qquad
 H=\{g\in G:g|_{D_\epsilon}=\operatorname{id}\}.
\]
Restriction $r:G\to E$ is locally trivial over its image. This is
Cerf's restriction theorem, including the boundary-incidence version
\cite[Theorem~1, p.~8-01, and Section~2.3, p.~8-05]{Cerf1962Fibration}.
Its local extension construction can be supported inside
$\operatorname{int}D$, so the outer boundary creates no extra
condition. We next verify directly that $E$ is path connected and
that
\[
 \pi_1(G,\operatorname{id})\longrightarrow\pi_1(E,j)
 \quad\hbox{is surjective}.                                  \tag{A.125.1}
\]
This also proves that the image of restriction is all of $E$, by
lifting successive pieces of a path from $j$ using these local
extensions.

For clarity, the finite-family shrinking argument needed here does
not require applying a compact-target statement to an open ball.
Let $(e_z)$ be a compact family of orientation-preserving smooth
embeddings of $D_\epsilon$ into $\operatorname{int}D$. Precomposition
with $x\mapsto ux$, from $u=1$ to a sufficiently small fixed $u>0$,
gives a homotopy through embeddings. Compactness supplies a positive
lower bound for the least singular values of $A_z=De_z(0)$, and
uniform continuity of the derivatives gives
\[
 \sup_{z,x\in D_\epsilon}\|De_z(ux)-A_z\|
 <\tfrac12\inf_z\sigma_{\min}(A_z)
\]
when $u$ is small enough. Consequently straight-line interpolation
from $e_z(ux)$ to $p_z+uA_zx$, where $p_z=e_z(0)$, remains injective
and immersive: after subtracting the invertible linear map $uA_z$,
the Lipschitz constant of the error is smaller than half its least
singular value. Its images remain in the unit ball, which is convex.
The affine endpoint itself lies there when $u$ is small, uniformly
in $z$. Shrinking $u$ further if necessary, the centers $p_z$ can be
moved linearly to zero, still with all images inside the unit ball.
The polar decomposition $A_z=Q_zP_z$ then allows the positive
matrices $P_z$ to be interpolated to the identity, while retaining
the factor $u$; the resulting embeddings are $x\mapsto uQ_zx$.
All steps are continuous in the $C^\infty$ topology.

For a single embedding, the connectedness of $SO(3)$ finishes a
path to $j$. For a loop based at $j$, the same construction has
basepoint path exactly the radial shrinking $j(x)\mapsto uj(x)$,
and after that keeps the basepoint $uj$ fixed. It changes the loop
to $x\mapsto uQ_zx$, with $Q_z$ a loop in $SO(3)$ based at the
identity. The actual rotations $Q_z:D\to D$ form a loop in $G$;
shrinking their restrictions uses the same basepoint path.
Thus the two basepoint transports cancel and prove (A.125.1).
This is the elementary ball version of Cerf's frame-space argument
\cite[Proposition~3, p.~8-10]{Cerf1962Fibration}.

Cerf's three-dimensional diffeomorphism theorem says that $G$ is
path connected; the exact version for $D^3$ is stated in
\cite[Chapter~I, Theorem~1.2, p.~8]{Laudenbach1974}.
The exact sequence of the restriction fibration now gives
\[
 \pi_1G\longrightarrow\pi_1E\longrightarrow\pi_0H
       \longrightarrow\pi_0G,
\]
so (A.125.1) and $\pi_0G=0$ imply $\pi_0H=0$.

We carefully relate this conclusion to the bottom-fixed annulus
$A=D\setminus\operatorname{int}D_\epsilon\cong S^2\times I$.
Restriction identifies $H$ with the annulus diffeomorphisms whose
infinite jets on the inner boundary agree with the identity.
It does not assert that they fix an open outer collar of that
boundary. Conversely, equality of all jets is exactly enough to
extend such an annulus diffeomorphism smoothly by the identity on
$D_\epsilon$; the inverse extends in the same way. These operations
identify the corresponding smooth mapping spaces and paths.

Every bottom-fixed smooth annulus diffeomorphism $h$ is isotopic,
through bottom-fixed diffeomorphisms, to one in this image. Indeed,
its upper restriction $a$ preserves orientation. The connectedness
of $\operatorname{Diff}^+(S^2)$ gives a smooth isotopy $a_s$ from
$\operatorname{id}$ to $a$; see
\cite[Theorem~4 and Corollary~2, p.~8-13]{Cerf1962Fibration}.
Choose a smooth function $\eta:I\to I$ equal to zero near zero and
equal to one near one. Postcomposing $h$ with
\[
 (x,t)\longmapsto(a_{s\eta(t)}^{-1}(x),t)
\]
changes it, within the bottom-fixed group, to a diffeomorphism $h'$
fixed on both boundary components. Cerf's boundary-jet proposition
\cite[Proposition~1 and its special case, p.~8-06]{Cerf1962Fibration},
applied with $M=S^2\times\{0\}$ and order $r=\infty$, gives the
needed normalization. More precisely, put
\[
 D_\infty=\operatorname{Diff}(A;J^\infty_{\partial A}),\quad
 D_M=\operatorname{Diff}(A;\partial A;J^\infty_M),\quad
 D_0=\operatorname{Diff}(A;\partial A).
\]
The proposition states that $D_\infty\to D_M$ induces
isomorphisms on $\pi_i$ for $i\geq0$; its displayed special case
states the same for the composite $D_\infty\to D_0$.
Consequently $D_M\to D_0$ induces an isomorphism on $\pi_0$
by two-out-of-three. Thus $h'$ is isotopic relative to the entire
boundary to an $h''\in D_M$. Extend $h''$ by the identity on $D_\epsilon$. Since $H$ is
path connected, restrict a path in $H$ to obtain a bottom-fixed
smooth isotopy of $h''$ to the identity. This proves
$\pi_0P_{\rm Diff}(S^2)=0$ without confusing pointwise fixed inner
disks, identity jets, and fixed collars.

Finally let $h$ be a PL concordance on $S^2$. The
three-dimensional category comparison
\cite[Chapter~I, Section~1, Corollary~3 and conclusion,
pp.~7--9]{Laudenbach1974} gives a TOP isotopy $J_s$ of the annulus
from $h$ to a smooth diffeomorphism $g$. No assertion that $g$ is
literally PL in the chosen coordinates is required. Every $J_s$
preserves each boundary component. If $b_s$ denotes its bottom
restriction, then
\[
 \widehat J_s=(b_s^{-1}\times\operatorname{id}_I)J_s
\]
is a TOP isotopy fixing the bottom. Its endpoint
$g'=(b_1^{-1}\times\operatorname{id}_I)g$ is smooth, since $b_1$
is the bottom restriction of the smooth map $g$.
We have proved that $g'$ is smoothly isotopic to the identity
through bottom-fixed maps. Thus $h$ and the identity are
TOP-isotopic. The injectivity of
$\pi_0\operatorname{Aut}_{\rm PL}(A)\to
 \pi_0\operatorname{Aut}_{\rm TOP}(A)$ in
\cite[Chapter~I, Section~1, pp.~7--8]{Laudenbach1974}
now gives an unrestricted jointly PL isotopy $L_s$ between the
two genuinely PL endpoints $h$ and the identity.
If $c_s$ is its bottom restriction, the total inverse boundary
family is PL, and
$(c_s^{-1}\times\operatorname{id}_I)L_s$ is a jointly PL
bottom-fixed isotopy. Its endpoints remain $h$ and the identity,
because $c_0=c_1=\operatorname{id}$. Concatenation proves the real
PL assertion. A.115 and the relative component argument of A.113
then give every real closed field and every definable model.
\end{proof}

% Independently audited and promoted 2026-11-19T17:51:45+00:00
\begin{saresult}{A.126}{Known}{Simply connected closed five-manifolds have trivial concordance components}
Let $M$ be a closed connected PL five-manifold with $\pi_1M=0$.
For any definable realization of its PL model over any real closed
field,
\[
 \pi_0P_{\rm PL}(M)=\pi_0P_{\rm def}(M)=0.
\]
There is no restriction on $\pi_2M$. The smooth input holds for
closed simply connected smooth manifolds in every dimension at
least five. Only the closed case of that input is invoked here.
\end{saresult}
\begin{proof}
The needed theorem is the full pseudo-isotopy theorem, not merely
the implication that pseudo-isotopic endpoint diffeomorphisms are
isotopic. In the original paper, Cerf defines the pseudo-isotopy
group on printed p.~7 as the diffeomorphisms of $V\times I$
fixed on $V\times\{0\}$. Theorem~0 on printed p.~8 states that
this group is connected for closed smooth $V$ with
$\pi_1V=0$ and $\dim V\geq5$
\cite[Introduction, Definition~2 and Theorem~0, pp.~7--8]{Cerf1970}.
The endpoint assertion occurs separately as Corollary~1 on that
page. The main connectivity theorem is proved as Theorem~3 on
printed p.~156, using the critical-point-free function-space
model described in the introduction.

The smooth group with the bottom fixed is locally path connected,
so this connectedness means that its component group is zero.
For completeness, choose a product metric near the boundary of
$V\times I$. Sufficiently $C^1$-small diffeomorphisms are expressed
uniquely as $x\mapsto\exp_x X(x)$ in the usual local-addition
chart. The boundary-preserving condition gives vector fields
tangent to the boundary, and the bottom-fixed condition gives
$X=0$ on the bottom. A small convex neighborhood of zero in this
linear space stays inside the diffeomorphism chart. Scaling $X$
to zero therefore gives paths through bottom-fixed smooth
diffeomorphisms. This proves the required local path connectedness
in the $C^\infty$ topology. Consequently Cerf's theorem gives
\[
 \pi_0P_{\rm Diff}(V)=0.                              \tag{A.126.1}
\]

Every PL five-manifold admits a compatible smoothing, without any
independently prescribed boundary smoothing. Here is the exact
low-dimensional input already used in A.124. In a finite PL handle
decomposition the index-$k$ extension obstruction is
$\Gamma_{k-1}$
\cite[Essay~IV, Appendix~B, p.~211]{KirbySiebenmann1977}.
The compatible low-dimensional smoothings in
\cite[Essay~I, p.~39]{KirbySiebenmann1977}, together with
$\Gamma_4=0$ explicitly on
\cite[Essay~IV, p.~200]{KirbySiebenmann1977}, remove all obstructions
for handles of index at most five. Let $M^{\rm sm}$ be a compatible
smoothing of $M$ and $\tau:|K|\to M^{\rm sm}$ its differentiable
triangulation. It is still closed and simply connected.

The precise category comparison now used is the component
isomorphism of Burghelea--Lashof
\cite[Theorem~6.1, p.~41, and proof pp.~42--43]{BurgheleaLashof1974II},
with the concordance groups defined on p.~41:
\[
 c_\tau:\pi_0P_{\rm Diff}(M^{\rm sm})
       \xrightarrow{\ \cong\ }\pi_0P_{\rm PL}(|K|).    \tag{A.126.2}
\]
Their bottom-and-side convention agrees with ours; the side is
empty here. This is the unstable component comparison associated
to the differentiable triangulation, as already used in A.103 and
A.124. It does not require a smooth map to be literally PL in a
fixed coordinate triangulation, and no stable comparison is being
substituted for (A.126.2). Equations (A.126.1)--(A.126.2) give
$\pi_0P_{\rm PL}(M)=0$.

Finally the finite PL field comparison A.115 and the relative
definable component comparison A.113 preserve this group and its
bottom condition. They prove the asserted conclusion over every
real closed field. The absence of a $\pi_2$ assumption is exactly
what extends the dimension-five range of A.124.
\end{proof}

% Independently audited and promoted 2026-11-19T17:51:45+00:00
\begin{saresult}{A.127}{Proved under stated hypotheses}{The projective plane has trivial concordance component group}
For the standard PL projective plane and every definable model over
any real closed field,
\[
 \pi_0P_{\rm PL}(\mathbb {RP}^2)
 =\pi_0P_{\rm def}(\mathbb {RP}^2)=0.
\]
The bottom is fixed pointwise and the top is free. Together with
A.123 and A.125 this computes the component group for every closed
surface. No higher parametric comparison is asserted.
\end{saresult}
\begin{proof}
Put $M=\mathbb {RP}^2\times I$. We work over the real field first.
Every homeomorphism and every isotopy constructed below is jointly
PL, including the twists used to change the annulus winding.

We use two elementary consequences of the three-dimensional
PL Schoenflies theorem, whose original sphere statement is
\cite[p.~6, with proof pp.~7--8]{Alexander1924}.
First, $M$ is irreducible. Its orientation double cover is
$S^2\times I$, with deck involution antipodal on $S^2$ and identity
on $I$. An embedded PL sphere in $\operatorname{int}M$ lifts to two
disjoint spheres exchanged by this involution. Cap the two ends of
$S^2\times I$ by balls to obtain $S^3$. For each lifted sphere,
Schoenflies gives the following dichotomy: either the two end caps
lie on the same side, and the other side is a ball lying wholly in
the open cylinder, or the sphere separates the two ends.
Two disjoint spheres separating the ends are strictly ordered from
bottom to top. An involution preserving each end cannot exchange
them. Hence the lifts bound interior balls. Balls with disjoint
boundaries in this cylinder are disjoint or nested; strict nesting
of the two balls is impossible because the involution has order two.
The balls are therefore disjoint, and either projects to a ball
bounded by the original sphere. This proves irreducibility.

Second, two proper PL disks in a PL three-ball with the same
boundary curve are ambient isotopic relative to the entire boundary
of the ball. To see this without an intersection-minimization
assumption, each disk separates the ball into two PL balls, by
Schoenflies applied to the disk joined to either boundary hemisphere.
Choose a PL homeomorphism between the disks, fixed on their common
boundary. On each complementary ball, glue this map to the identity
on the corresponding boundary hemisphere. The resulting boundary
sphere homeomorphism extends across that ball by coning. The two
extensions glue to a homeomorphism of the original ball fixed on
its entire boundary and taking the first disk to the second. The
Alexander isotopy contracts this ball homeomorphism relative to its
boundary, proving the assertion. The same conical formula shows
that this isotopy is jointly PL, as in A.123.

We describe the standard annuli and suitable PL twists explicitly.
Let $D$ be a centrally symmetric polygonal disk, identify its
boundary with $\mathbb R/2\mathbb Z$ by an edgewise linear
coordinate $s$, and make the identification compatible with
$x\mapsto-x$ by $s\mapsto s+1$. The quotient of $D$ by antipodal
identification on its boundary is $\mathbb {RP}^2$.
The image of $\partial D$ is a one-sided projective line $c$.
The image of a diameter of $D$ joining the points $s=0$ and $s=1$
is another one-sided projective line $a$, meeting $c$ once. Put
\[
 C=c\times I,\qquad A_0=a\times I.
\]
Both are proper one-sided annuli. Cutting $M$ along either one
gives $D^2\times I$, a three-ball. In particular cutting along $C$
is exactly the quotient description $D\times I\to M$.

We first construct a jointly PL disk isotopy from the identity to
the exact linear antipodal map $J:x\mapsto-x$, such that every
boundary map commutes with $J$. The boundary translation
$c_u(s)=s+u$ on $\mathbb R/2\mathbb Z$ is jointly PL. On a PL
collar $\partial D\times[0,\delta]$, with $v=0$ the outer boundary,
extend it by
\[
 (s,v)\longmapsto
 \bigl(s+\max\{0,u-v/\delta\},v\bigr),
 \qquad 0\leq u\leq1,
\]
and use the identity outside the collar. This gives a jointly PL
isotopy $r_u$ of $D$ whose final boundary map is $J$. The map
$b=J^{-1}r_1$ fixes $\partial D$ pointwise. By the relative disk
Alexander isotopy, join $b$ to the identity through
boundary-fixed PL maps $b_u$. Concatenating $r_u$ with $Jb_u$
gives a jointly PL isotopy from the identity to exactly $J$.
On the first half its boundary map is translation by a number from
zero to one; on the second half it is constantly translation by
one. In particular every boundary map commutes with $s\mapsto s+1$.

Repeat this path, composing successive copies with $J^j$, to obtain
for each positive integer $k$ a jointly PL path $q_t$ from the
identity to $J^k$. For negative $k$ use the inverse path on each
copy; for $k=0$ use the constant path. The lift of its boundary-translation amount starts at zero
and ends at $k$, and is constant only on the Alexander-correction
parts. Since all boundary maps commute with the antipodal
identification, $q_t$ descends to a jointly PL isotopy $\bar q_t$
of $\mathbb {RP}^2$. Define
\[
 T_k(x,t)=(\bar q_t(x),t),\qquad
 T_{k,u}(x,t)=(\bar q_{\min\{u,t\}}(x),t).             \tag{A.127.1}
\]
The second family is a jointly PL bottom-fixed isotopy from the
identity to $T_k$. Set $A_k=T_k(A_0)$. Its intersection with $C$
is the graph of the boundary-translation amount modulo one, hence
a proper arc of winding number $k$ with the same endpoints as
$A_0\cap C$. Its bottom and top boundary curves are both $a$,
since $q_0$ is the identity and $q_1=J^k$ preserves the diameter.
This supplies actual PL twists and their PL contractions; no polar
rotation formula is being called PL.

Let $h\in P_{\rm PL}(\mathbb {RP}^2)$. Its upper restriction is
homotopic to the identity through
$x\mapsto\operatorname{pr}_1h(x,t)$. The surface homotopy-isotopy
theorem \cite[Theorem~6.3, p.~99]{Epstein1966} gives a PL isotopy
$a_u$ of the upper map $a$ to the identity. Put $c_u=a_ua^{-1}$
and extend this inverse correction by
$C_u(x,t)=(c_{\min\{u,t\}}(x),t)$. This is a jointly PL
bottom-fixed isotopy starting at the identity, and $C_1h$ has
upper restriction $a^{-1}a=\operatorname{id}$.
After this bottom-fixed PL isotopy we may assume that
$h$ fixes both boundary components pointwise.

Put $A=h(A_0)$. Arrange $A$ and $C$ in PL general position
relative to the boundary. Their boundaries meet in exactly two
points, one on each end. Thus their intersection consists of one
proper spanning arc and finitely many circles. Every such circle
is inessential in each annulus: an essential circle in an annulus
must meet every arc joining its two boundary components, whereas
the intersection components are disjoint. Choose a circle innermost
on $C$. It bounds a disk $D_C\subset C$ with interior disjoint
from $A$, and also a disk $D_A\subset A$. Their union is an
embedded PL sphere in the interior of $M$, and bounds a ball $B$
by irreducibility. The surface
$(A\setminus\operatorname{int}D_A)\setminus\partial D_A$ is
connected, contains the two original boundary circles of $A$ on
$\partial M$, and is disjoint from this sphere. It consequently
lies outside $B$. In particular the
interior of $B$ misses $A$. Pushing $D_A$ across this ball and
slightly past $D_C$ removes the selected circle and any circles
inside $D_A$, introduces no new intersection components, and is
supported away from $\partial M$. Repeating finitely often leaves
a single spanning arc.

This remaining arc has an integer winding number $k$ in $C$,
relative to its fixed endpoints. The relative arc theorem
\cite[Theorem~3.1, p.~89]{Epstein1966} isotopes it to
$A_k\cap C$, keeping $\partial C$ fixed. Extend this isotopy to
the ambient manifold, relative to $\partial M$, using the normal
interval-bundle construction described below. After this change,
cut along $C$. The annuli $A$ and $A_k$ each cut to a proper PL
disk in the ball $D\times I$, and their boundary curves agree:
their two side arcs are the two lifts of $A_k\cap C$, and their
end arcs are the common diameter on the top and bottom.
The disk assertion already proved gives an ambient PL isotopy
fixed on the entire cut-ball boundary taking one disk to the other.
It descends to $M$. Therefore an ambient PL isotopy relative to
$\partial M$ has changed $h$ to a map $h_1$ with
$h_1(A_0)=A_k$.

Compose with $T_k^{-1}$ and put $g=T_k^{-1}h_1$. This is still PL,
fixes the bottom, and preserves $A_0$ setwise. It is connected to
$h$ by a bottom-fixed PL isotopy, by the explicit path for $T_k$.
Its restriction $f=g|_{A_0}$ is a bottom-fixed annulus
homeomorphism. The circle calculation in A.123 gives a jointly PL
isotopy $f_u$ from $f$ to the identity, fixed on the bottom circle.
Set $c_u=f_uf^{-1}$, an isotopy from the identity to $f^{-1}$.

Here is the promised extension, including the one-sided case.
A boundary-compatible normal neighborhood of the proper annulus
$A_0$ is
\[
 N=(\widetilde A_0\times[-\delta,\delta])/
       ((p,v)\sim(dp,-v)),
\]
where $\widetilde A_0$ is its normal-orientation double cover and
$d$ is the deck involution. Lift $c_u$ starting at the identity to
$\widetilde c_u$. The lifts commute with $d$ and fix every lift
of a boundary point fixed by $c_u$. Define
\[
 B_u([p,v])=
 \left[\widetilde c_{\tau(u,|v|)}(p),v\right],\qquad
 \tau(u,r)=\max\{0,u-r/\delta\}.                    \tag{A.127.2}
\]
This is well defined, is the identity for $|v|=\delta$, and has
inverse given by the same formula with the inverse lifted family.
It therefore extends by the identity outside $N$ to an ambient
isotopy. Covering lifts and composition of the finite PL maps show
that it is jointly PL. It fixes the whole bottom because $c_u$
fixes its bottom circle. For the preceding arc isotopy in $C$ the
same construction fixes both boundary components, as that isotopy
fixes both boundary circles. Thus no two-sided normal framing was
assumed at either use.

The map $g'=B_1g$ now fixes $A_0$ pointwise. Its action on the two
normal directions is locally constant along the connected annulus.
It is the identity at the bottom, since the entire bottom surface
is fixed. Hence it preserves the normal directions everywhere.
Cutting along $A_0$ consequently lifts $g'$ to a PL homeomorphism
of $D^2\times I$ fixing its bottom and its whole side pointwise.
The ball calculation of A.123 gives an isotopy to the identity
relative to these sets: first cancel the free top-disk map by its
boundary-relative Alexander isotopy and a collar extension, and
then Alexander-contract the resulting full-boundary-fixed
three-ball map. This isotopy descends to $M$. We have proved that
the original PL map $h$ is jointly PL isotopic to the identity through
bottom-fixed maps. This proves the real PL assertion. A.115 and the
relative component argument of A.113 give the assertion over every
real closed field and for every definable model.

The integer $k$ was used only to select an arc model after a chosen
intersection simplification. We neither asserted that it is an
ambient-isotopy invariant nor that $T_1$ fixes the upper projective
plane. Those stronger assertions are unnecessary and false in the
literal rotation model.
For the precise source correction, the displayed twist in
\cite[Appendix~A, Theorem~A.6, pp.~614--615]{BoyleChen2024}
sends the upper interior point $(1/2,0)$ of the disk model to
$(-1/2,0)$. They are distinct quotient points, because only
boundary points are identified. Thus that literal twist does not
belong to the full-boundary-fixed group claimed there. The exact
two-point check is recorded with the source audit; the assertion of
that theorem is not an input to the present free-top proof.
\end{proof}


\begin{saresult}{A.128}{Known}{The three- and four-sphere component cases follow directly in PL}
For every standard PL sphere $S^n$, $n\geq2$,
\[
 \pi_0 P_{\mathrm{PL}}(S^n)=0.
\]
Consequently the same assertion holds for each definable model of
the standard PL sphere over every real closed field, with the
simplicial concordance convention of this study. In particular,
this closes the standard $S^3$ and $S^4$ component cases without
a smooth-to-PL comparison in base dimension three or four.
\end{saresult}
\begin{proof}
Write
\[
 X_j=\mathrm{PL}_j/\mathrm O_j,
\]
using the simplicial groups of germs and their homogeneous spaces.
These are the unstable smoothing spaces, not their stable
substitutes. The sphere calculation in
\cite[\S6, p.~43, equation~(1), following Proposition~5.3]{BurgheleaLashof1974II}
gives a weak equivalence of spaces
\[
 P_{\mathrm{PL}}(S^n)
 \simeq E(X_{n+1},X_n)
 :=\operatorname{hofib}_{*}(X_n\longrightarrow X_{n+1}).
                                                        \tag{A.128.1}
\]
Here $E(X_{n+1},X_n)$ is the path space defined in the same
paper, p.~42: a point consists of a point of $X_n$ and a path
in $X_{n+1}$ from the basepoint to its image. The sphere calculation
uses the Kuiper--Lashof fibration, which the paper states on
p.~38 before Proposition~5.3 as
\[
 \mathrm{PL}_n\longrightarrow\mathrm{PL}_{n+1}
 \longrightarrow S^n\times B P_{\mathrm{PL}}(S^n),
\]
with the orthogonal sphere fibration mapping into it. It has no
high-dimensional stability hypothesis. The original
\cite[Theorem~1 and definitions, pp.~243--244]{KuiperLashof1966II}
defines each simplex by a PL bundle isomorphism on the entire
product $\Delta_k\times X$, with base $\Delta_k$; its concordance
group fixes $S^n\times0$ pointwise and leaves the top free.
On these compact products the PL maps have finite triangulations,
so the convention is precisely the jointly PL convention used here.
Only path components of
(A.128.1) will be used, so no compatibility of this equivalence
with a proposed multiplication is required.

The low-degree stabilization statement
\cite[\S5, p.~37, formula~(3)]{BurgheleaLashof1974II}
is
\[
 \pi_i(X_j)\xrightarrow{\ \cong\ }
 \pi_i(\mathrm{PL}/\mathrm O)\qquad(i<j).
                                                        \tag{A.128.2}
\]
The stable space $\mathrm{PL}/\mathrm O$ is connected and has
vanishing homotopy groups through degree six; the relevant
vanishing is explicitly recorded in
\cite[Annex~C, p.~318, calculation following Theorem~5.4]{KirbySiebenmann1977}.
For $n\geq2$, (A.128.2) therefore gives
\[
 \pi_0(X_n)=0,\qquad \pi_1(X_{n+1})=0.
\]
The exact homotopy sequence of (A.128.1) now contains
\[
 \pi_1(X_{n+1})\longrightarrow
 \pi_0 P_{\mathrm{PL}}(S^n)\longrightarrow\pi_0(X_n),
\]
and both outer terms are trivial. The middle pointed set has
one element, hence the concordance component group is trivial.

For a general real closed field, A.115 transports the real PL
component computation. The relative one-parameter component
comparison in A.113 transports it to the stipulated definable
group. A semialgebraic sphere may be replaced by its finite PL
model before this comparison.

There are two distinct statements in the source that must not be
confused here. Formula~(3) on p.~37 is the low-degree assertion
$i<j$ used above. The stronger Theorem~5.1 on that page is a
high-dimensional stability assertion and is not used. Likewise,
Theorem~6.1 on p.~41, whose printed statement omits a dimension
restriction although its displayed proof invokes Theorem~5.1,
is not used in base dimension three or four. The argument gives
no component computation for a general closed three- or
four-manifold.
\end{proof}


\begin{saresult}{A.129}{Known}{All closed simply connected three-manifolds have trivial concordance components}
Let $M$ be a closed connected simply connected PL three-manifold. Then
\[
  \pi_0P_{\mathrm{PL}}(M)=0.
\]
The same holds for the stipulated definable concordance group of every
closed definably simply connected definable three-manifold over a real
closed field.
\end{saresult}
\begin{proof}
Give $M$ a compatible smooth structure, whose existence in dimensions
at most four is explicitly recorded in
\cite[Essay~I, p.~39]{KirbySiebenmann1977}.
Simple connectivity implies orientability. Perelman's elliptization
conclusion \cite[Section~1.4, p.~3]{Perelman2003Extinction} identifies
this smooth manifold with a spherical space form $S^3/\Gamma$.
Since $S^3$ is its universal cover, $\Gamma$ is its fundamental group;
therefore $\Gamma=1$. In particular there is a topological homeomorphism
$h:M\longrightarrow S^3$.

To keep track of the given PL structures, apply the original
approximation theorem in
\cite[Approximation Theorem, p.~34; proof Section~9, pp.~75--79]{Shalen1984}
to $h$. It gives a PL embedding $q:M\longrightarrow S^3$.
Both boundaries are empty, so the relative-boundary hypotheses are
vacuous. The image of $q$ is open by invariance of domain and closed
by compactness. Since $S^3$ is connected and the image is nonempty,
$q$ is onto and is a PL homeomorphism. This use of Shalen's proved
refinement of the three-dimensional approximation theorem does not
presuppose that Moise's 1952 original has been opened.

Conjugation by $q\times\operatorname{id}_I$, and by its product with
every parameter simplex, is an isomorphism of the jointly PL
concordance simplicial groups
\[
  P_{\mathrm{PL}}(M)\cong P_{\mathrm{PL}}(S^3).
\]
A.128 makes its component group trivial. Finally, take a finite PL
model for the definable manifold and transfer it to the real field.
Realize its finite abstract triangulation with rational vertices,
choosing a rational vertex as basepoint. This is a pointed
semialgebraic set defined without parameters. The natural comparison
\cite[arXiv version~1, Corollary~4.3, p.~10]{BaroOtero2008v1}
identifies its o-minimal fundamental group over the original field
with the classical fundamental group of its real realization. Thus
the real model is simply connected. This citation uses the explicit
preprint numbering, not the numbering of the final published paper.
A.115 and the relative component comparison A.113 then
transfer the result back to the original field and expansion.
No smooth-to-PL concordance comparison in base dimension three is
used here.
\end{proof}

\begin{saresult}{A.130}{Known}{Infinite-rank concordance subgroups in closed orientable dimension three}
Let $B$ be a closed connected orientable PL three-manifold with
infinite fundamental group. There is a group monomorphism
\[
 \bigoplus_{j\ge1}\mathbb Z\ \hookrightarrow
 \pi_0P_{\mathrm{PL}}(B).
                                                        \tag{A.130.1}
\]
The same holds for its definable realizations over every real
closed field. Over the real field, these classes remain independent
after forgetting to topological concordances, even after further
forgetting the fixed-bottom condition.
This is a subgroup result, not a presentation of the entire
three-dimensional concordance component group.
\end{saresult}
\begin{proof}
Give $B$ a compatible smooth structure by
\cite[Essay~I, p.~39]{KirbySiebenmann1977}; put $Z=B\times I$ and
$K=B\times0$. There are no corners because $B$ is closed.
Choose a differentiable triangulation
$t:|Z|\longrightarrow Z$ compatible with $K$ and with the given
PL product structure.

We first construct the comparison on components with all maps typed.
Let $P_{\mathrm{PD}}(Z,t,K)$ be the semisimplicial complex of
piecewise differentiable homeomorphisms from the triangulated
$|Z|$ to $Z$, with $t$ prescribed on $K$, and with total parameter
maps preserving each parameter. It need not be a group.
There are maps of semisimplicial spaces
\[
 \begin{aligned}
 j:P_{\mathrm{PL}}(B)&\longrightarrow P_{\mathrm{PD}}(Z,t,K),
       &j(p)&=t\circ p,\\
 i:P_{\mathrm{Diff}}(B)&\longrightarrow P_{\mathrm{PD}}(Z,t,K),
       &i(h)&=h\circ t.
 \end{aligned}
                                                        \tag{A.130.2}
\]
The original relative PL--PD comparison makes $j$ a homotopy
equivalence, without a dimension restriction:
\cite[definitions and comparison, pp.~1--2; Theorem~1.1(c), p.~8]{BurgheleaLashof1974I}.
Thus
\[
 c=(\pi_0j)^{-1}\pi_0i:
 \pi_0P_{\mathrm{Diff}}(B)\longrightarrow
 \pi_0P_{\mathrm{PL}}(B)
                                                        \tag{A.130.3}
\]
is well-defined as a map of sets.

It is a group homomorphism. If $h_at$ is joined through PD
homeomorphisms to $tp_a$, for $a=1,2$, then
\[
 h_1h_2t\ \simeq\ h_1tp_2\ \simeq\ tp_1p_2.
\]
The first homotopy uses the left action of smooth homeomorphisms
on PD maps; the second uses the right action of PL homeomorphisms.
Both preserve $K$ and all parameters. Injectivity of $\pi_0j$
therefore proves $c([h_1h_2])=c([h_1])c([h_2])$.
Forgetting PD structure turns each of these paths into a TOP
isotopy. Consequently the composite of $c$ with the PL-to-TOP
map on components is the ordinary smooth-to-TOP map.
No injectivity or surjectivity of $c$ is asserted in dimension three.

The source comparison restricts to maps homotopic to the identity
as maps of the boundary pair, relative to $K$. Every pseudo-isotopy
satisfies this condition. Indeed, write
$h(x,u)=(a(x,u),b(x,u))$. First deform its second coordinate
linearly from $b(x,u)$ to $u$. Then deform
$(a(x,u),u)$ to $(x,u)$ through
\[
 (x,u)\longmapsto(a(x,(1-s)u),u),\qquad 0\le s\le1.
\]
These are continuous maps of the boundary pair, fixed on $K$;
they need not be homeomorphisms, as the convention asks only for
a homotopy of maps. Hence no pseudo-isotopy class is removed by
that convention.

We now use a stronger form of the detection behind
\cite[Theorem~B, p.~2; Theorem~C, p.~13 and its deduction on p.~14;
Theorem~D, p.~28]{LinXieZhang2026}.
Those results show that the image of
\[
 \pi_0\operatorname{Diff}_{\partial}(I\times B)
 \longrightarrow \pi_0\operatorname{Homeo}(I\times B)
                                                        \tag{A.130.4}
\]
contains an abelian subgroup of infinite rank. In the irreducible
case this is Theorem~C; in the reducible case it is Theorem~D.
Both use diffeomorphisms equal to the identity near the entire
boundary. Their domain component group is abelian, as noted
on p.~1 of that source. Choose a countable family of classes
in that domain whose images in (A.130.4) are linearly independent
over $\mathbb Z$; such a family exists by infinite rank.

Include them in $\pi_0P_{\mathrm{Diff}}(B)$ and apply $c$.
Their images commute because the source classes commute.
A finite integral relation between their images in
$\pi_0P_{\mathrm{PL}}(B)$ would, by the commuting TOP triangle,
give the same relation in (A.130.4). The chosen independence
forces every coefficient to vanish. This proves (A.130.1).
Detection in the completely boundary-free homeomorphism group
also shows why no comparison between a fixed boundary collar
and a merely pointwise-fixed bottom is required here.

For other real closed fields, transport the PL result using
A.115 and then apply the relative component comparison A.113.
The assertion for a definably orientable three-manifold with
infinite definable fundamental group follows by choosing a
rational finite model and using the exact homotopy comparison
\cite[arXiv version~1, Corollary~4.3, p.~10]{BaroOtero2008v1}.
Nothing in this argument evaluates the full component group,
uses the low-dimensional extension of Burghelea--Lashof~II,
or settles the higher-parametric comparison.
\end{proof}

\clearpage
\begin{saresult}{A.131}{Proved under stated hypotheses}{Effective countable models and complete recursive presentations for PL and semialgebraic concordances}
Let $K$ be a nonempty finite abstract simplicial complex whose realization
is a closed PL manifold. The ordered finite list of simplices of $K$ is
all the input in the following construction. There are explicit
countable simplicial groups
\[
 G^{\mathrm{PL}}(K),\qquad G^{\mathrm{SA}}(K)
                                                        \tag{A.131.1}
\]
with the following effective properties. Their simplices are finite
integer strings; membership, equality, and degeneracy are decidable.
Multiplication, inversion, faces and degeneracies are given by total
algorithms. The construction also gives a computable simplicial
homomorphism $G^{\mathrm{PL}}(K)\to G^{\mathrm{SA}}(K)$.
For every real closed field $R$, coefficient extension gives weak
equivalences
\[
 G^{\mathrm{PL}}(K)\longrightarrow P_{\mathrm{PL},R}(|K|_R),
 \qquad
 G^{\mathrm{SA}}(K)\longrightarrow P_{\mathrm{SA},R}(|K|_R),
                                                        \tag{A.131.2}
\]
where the source is realized over the real algebraic numbers.
All families are homeomorphisms of the total product with the
parameter simplex, preserve that entire parameter, fix the bottom
pointwise, and preserve the free top setwise.

For each $r\ge0$, the algorithm in the proof outputs a complete
countable presentation $\Pi_r^{\mathrm{PL}}(K)$, respectively
$\Pi_r^{\mathrm{SA}}(K)$, for the entire homotopy group
\[
 \pi_rP_{\mathrm{PL},R}(|K|_R),\qquad
 \pi_rP_{\mathrm{SA},R}(|K|_R).
                                                        \tag{A.131.3}
\]
For $r\ge1$ these groups are based at the identity; translation
identifies the corresponding groups on any other component.
Here ``outputs'' means that one uniform algorithm enumerates all
specified generators and relators. For $r\ge1$ it may instead
enumerate an integer relation matrix for the abelian group.
For an arbitrary fixed o-minimal expansion, $\Pi_0^{\mathrm{PL}}(K)$
also presents the stipulated definable component group. If every
component of $|K|$ has dimension at least five, (A.131.1) gives a
weak model for the real TOP concordance space as well, and
$\Pi_r^{\mathrm{PL}}(K)$ presents its homotopy groups for all $r$.

For the empty complex and empty base, take the constant one-element
simplicial group; every component and homotopy group is trivial.

These are complete effective countable presentations and whole
simplicial models. No terminating word problem, finite presentation,
finite invariant-factor description, or efficient complexity bound
is asserted. No higher weak equivalence from PL to SA or to an
arbitrary o-minimal expansion is inferred.
\end{saresult}
\begin{proof}
\emph{The finite coordinate model and its arithmetic.}
Order the $v$ vertices of $K$ and realize them as the standard basis
vectors of $\mathbb Q^v$. Put
\[
 B=|K|,\quad Z=B\times[0,1],\quad A=B\times\{0\},\quad
 T=B\times\{1\},\quad D_q=Z\times\Delta^q.
\]
Membership in $B$ is the finite disjunction, over the maximal
simplices $\sigma$ of $K$, of
\[
 x_i\ge0,\qquad \sum_i x_i=1,\qquad x_i=0\ (i\notin\sigma).
\]
This is a fixed rational linear formula. The other displayed sets have
rational linear formulas too. None of the algorithms needs to
recognize whether the promised finite complex is a manifold.

Write $F=\mathbb Q^{\mathrm{rc}}$ for the ordered real algebraic
numbers. One may encode an element by a nonzero integer polynomial,
a rational open interval with exactly one root, and that root.
The final phrase refers to its unique value, not additional input:
the finite input is the polynomial and the two rational endpoints.
Sturm root counting decides whether this code is valid. Every real
algebraic number has such a code, and the set of valid codes is
effectively enumerable. Arithmetic, sign and equality can be decided
by root isolation, or by the following uniform reduction to rational
quantifier elimination. Replace each coded number by a fresh variable,
constrain it by its polynomial and isolating interval, and ask the
corresponding polynomial sentence. Each such variable has exactly one
possible value, so this replacement preserves truth. An output code
for an arithmetic operation is found by enumerating root codes and
deciding equality to the desired expression; a code exists, so the
search terminates. Duplicate codes for a number cause no difficulty.

The exact quantifier-elimination input is
\cite[online edition dated 26 June 2016, Theorem~2.94, p.~80;
Algorithm~14.16, pp.~586--587; Algorithm~14.21, pp.~591--592,
and Theorem~14.22, pp.~592--593]{BasuPollackRoy2016Online}.
The same edition states Sturm's theorem as Theorem~2.62, p.~62.
Thus the arithmetic used below is an algorithm on finite integers,
not arithmetic with an oracle for arbitrary real constants.
For linear formulas over this effective ordered field there is also
an elementary elimination algorithm: put a Boolean formula into
disjunctive normal form; substitute an equality with nonzero
coefficient when present; otherwise compare each lower bound with
each upper bound for the variable being eliminated. Retain a strict
comparison if either original bound was strict. A finite family of
one-sided bounds is satisfiable when there is no opposite bound,
and constant inequalities are decided by their signs. Negating
converts universal to existential quantifiers. Iterating eliminates
all variables. This proves the linear decision claim used below.

\emph{PL graphs and their finite validity test.}
In degree $q$ use ambient coordinates $(z,z',t)$, where
$z,z'\in F^{v+1}$ and $t\in F^{q+1}$. In particular the source
and target share the \emph{same} parameter variable $t$.
A raw code is a finite list of closed simplices with vertices in $F$
in this ambient space. Affine independence is decided by Gaussian
elimination over $F$; reject a purported simplex that is not
independent, and allow zero-dimensional simplices. The list need not
already form a simplicial complex. Its union $\Gamma$ has the formula
\[
 \bigvee_{\sigma}
 \left\{w=\sum_{j=0}^{a_\sigma}\lambda_jv_{\sigma j},\quad
            \lambda_j\ge0,\quad\sum_j\lambda_j=1\right\}.
                                                        \tag{A.131.4}
\]
The braces abbreviate existential quantification of the finitely
many $\lambda_j$. All coefficients are constants in $F$, so
(A.131.4) is a linear formula.

Call the raw code valid precisely when these five finite sentences
hold:
\begin{enumerate}
\item $\Gamma\subset Z\times Z\times\Delta^q$;
\item for every $(z,t)\in D_q$ there exists exactly one $z'$ with
$(z,z',t)\in\Gamma$;
\item for every $(z',t)\in D_q$ there exists exactly one $z$ with
$(z,z',t)\in\Gamma$;
\item whenever $(z,z',t)\in\Gamma$ and $z\in A$, one has $z'=z$;
\item whenever $(z,z',t)\in\Gamma$, one has $z\in T$ if and only if
$z'\in T$.
\end{enumerate}
Uniqueness is the formula with two candidate outputs implying their
equality. Every sentence is linear and therefore decidable by the
elimination just given. Equality of two graph codes is the linear
sentence asserting their membership predicates equivalent at every
point, so it is decidable too.

We prove that this test is sufficient, including continuity. The
projection of $\Gamma$ to $(z,t)$ is bijective by the first three
conditions. On each simplex $\sigma$ it is affine and injective.
An affine map injective on a simplex is injective on its affine hull:
a kernel direction would give two distinct nearby points in the
relative interior with the same image. Its inverse on its image is
therefore affine. It follows that the function determined by
$\Gamma$ is affine on each of the finitely many projected simplices.
Those simplices are closed in $D_q$, and their affine formulas agree
on their overlaps because the graph is single valued. The finite
closed-cover gluing lemma proves continuity over $F$. Applying the
same argument to the other projection proves inverse continuity.

For completeness, a finite polyhedral subdivision compatible with
all the projected simplices is obtained by cutting with their affine-hull equations and ambient
extensions of all their face supporting hyperplanes, taking every
resulting face in every relevant affine span, and
triangulating those cells inductively by coning from an interior
point. A bounded nonempty polyhedron over $F$ has vertices over $F$,
by Gaussian elimination; the average of the vertices lies in its
relative interior. Thus all these subdivisions have vertices in $F$.
The affine formulas are compatible on this subdivision. The function
and its inverse are jointly PL, and the final two sentences impose
exactly the stipulated boundary conditions. Here PL means finite piecewise linear, as in A.115. Conversely, a jointly PL
homeomorphism on the finite polyhedron $D_q$ is affine on a finite
subdivision; its graph on each simplex is a simplex of the prescribed
kind. Hence every stipulated PL family has a valid graph code.
This equivalence used neither a homeomorphism-recognition oracle nor
ordinary compactness of a non-Archimedean field.

\emph{Canonical representatives and total operations.}
Encode all finite syntax above by binary strings using fixed
self-delimiting encodings of signed integers, lists and root codes.
Order the strings first by length, then lexicographically. There
are only finitely many predecessors of a given string. The canonical
representative of a valid graph is its least valid raw code in that
order. Given a valid code $c$, check the finitely many strings at
most $c$ for validity and graph equality, and take the first match.
Thus canonicalization terminates, and the predicate saying that a
string is canonical is decidable. Define $G_q^{\mathrm{PL}}(K)$ to
be the set of canonical valid degree-$q$ codes.

Here are finite formulas for all operations. If $g$ and $h$ are
represented by graph predicates $\Gamma_g$ and $\Gamma_h$, put
\[
 \begin{split}
 \Gamma_{g^{-1}}(z,z',t)&=\Gamma_g(z',z,t),\\
 \Gamma_{gh}(z,z',t)&=\exists y\,
       (\Gamma_h(z,y,t)\wedge\Gamma_g(y,z',t)).
 \end{split}
                                                        \tag{A.131.5}
\]
The product convention is $g$ after $h$.
For an order-preserving map $\alpha:[m]\to[q]$, let
\[
 (\alpha_*s)_j=\sum_{\alpha(i)=j}s_i.
\]
The simplicial pullback has graph
\[
 \Gamma_{\alpha^*g}(z,z',s)=
     [\,(z,s)\in D_m,\ (z',s)\in D_m\,]
       \wedge\Gamma_g(z,z',\alpha_*s).
                                                        \tag{A.131.6}
\]
This covers degeneracies as well as faces. The identity graph is
$z'=z$ on $D_q$.

All these are linear formulas. They are graphs of jointly PL
homeomorphisms: inversion exchanges the affine subdivisions;
composition uses intersections of the image cells of $h$ with the
domain cells of $g$; pullback uses the inverse images of the cells
under the affine map $(z,s)\mapsto(z,\alpha_*s)$ and the formula
for the first coordinate of $g$. Their inverse maps are obtained by
the same construction. These are finite bounded polyhedral covers,
triangulable over $F$ by the construction above. Consequently each
operation formula has at least one valid raw graph code.

The following uniform search outputs its canonical code: enumerate
raw codes in the prescribed order, reject invalid codes, and test
whether their graph predicate equals the operation formula. Equality
is decidable and a successful code exists, so the search terminates.
This proves total computability of every operation without leaving
a triangulation subroutine unspecified. Canonical equality is actual
equality of graphs, so associativity, inverse laws, and all
simplicial identities hold strictly. We have constructed a simplicial
group literally isomorphic to the jointly PL concordance group over
$F$.

\emph{The semialgebraic construction.}
For $G^{\mathrm{SA}}$ replace a list of graph simplices by a finite
Boolean formula in polynomial equations and inequalities with coded
coefficients in $F$, in the same variables $(z,z',t)$. Use the same
five validity conditions. Add joint continuity of the graph function
and of its inverse. These are first-order sentences, not assertions
that each individual fibre map is continuous. For example continuity
of the function is the sentence
\[
 \begin{split}
 \forall (z,z',t)\in\Gamma\ \forall\epsilon>0\ \exists\delta>0\
 \forall (w,w',s)\in\Gamma:
 \quad &\|w-z\|^2+\|s-t\|^2<\delta^2\\
 &\Longrightarrow
 \|w'-z'\|^2+\|s-t\|^2<\epsilon^2.
 \end{split}
                                                        \tag{A.131.7}
\]
Here squared norms are finite sums of coordinate squares, so the
formula is polynomial. Exchange the two graph projections to impose
inverse continuity. These conditions express precisely that the
total-space bijection is a homeomorphism. Quantifier elimination
decides their truth and equality of graphs. It also eliminates the
quantifiers in (A.131.5) when composition is formed. The same
least-string procedure produces canonical codes and total
algorithms. It follows exactly as before that
$G^{\mathrm{SA}}(K)$ is a simplicial group, literally the stipulated
semialgebraic concordance group over $F$.

Formula (A.131.4) is already a semialgebraic graph formula. Canonicalize
it in the SA coding to obtain the indicated computable map
$G^{\mathrm{PL}}\to G^{\mathrm{SA}}$. The map respects every
operation because each operation was defined by its actual graph.
A.115 and A.116, applied to $F\subset R$, prove (A.131.2).
The unique ordered embedding of the real algebraic field into any
real closed field specifies those coefficient extensions. There is
no map from non-Archimedean points to ordinary real points in this
argument.

These are whole effective simplicial models, not merely lists of
candidate components. For example, a degree-$q$ canonical simplex
$a$ is degenerate exactly when
\[
 a=s_i d_i a\quad\hbox{for some }0\le i<q.
\]
The finite test is decidable. Thus one can enumerate the
nondegenerate cells of the geometric realization and compute each
of their face attachments. This specifies its entire countable CW
model, with all incidence maps, in addition to its simplicial group
structure.

\emph{All homotopy groups with every relation.}
Let $G$ denote either effective simplicial group. Put
\[
 Z_0=G_0,\qquad
 Z_r=\bigcap_{i=0}^{r}\ker(d_i:G_r\to G_{r-1})\quad(r\ge1),
 \qquad
 N_{r+1}=\bigcap_{i=1}^{r+1}\ker(d_i:G_{r+1}\to G_r).
                                                        \tag{A.131.8}
\]
Membership in these sets is decidable, since there are finitely many
computable face tests. The simplicial identities show
$d_0N_{r+1}\subset Z_r$. This is a subgroup, and it is normal in
$Z_r$: for $a\in Z_r$ and $u\in N_{r+1}$,
\[
 (s_0a)u(s_0a)^{-1}\in N_{r+1},\qquad
 d_0\big((s_0a)u(s_0a)^{-1}\big)=a(d_0u)a^{-1}.
                                                        \tag{A.131.9}
\]
The same argument works in degree zero with $Z_0=G_0$.

Every simplicial group is Kan, and its Moore quotient computes its
homotopy groups. The precise opened source is
\cite[Theorem~17.1, pp.~67--68; Proposition~17.2, p.~68;
the second Proposition~17.3, pp.~68--69; Proposition~17.4, p.~69;
Theorem~16.6, p.~65]{May1967Simplicial}.
The original prints two propositions numbered 17.3; the second is
the normal-subgroup statement. It uses the last face convention,
whereas (A.131.8) uses the first. Reversing the order on each
standard simplex exchanges these conventions. More directly, if
$z$ represents an elementary simplicial homotopy with
\[
 d_0z=x,\quad d_1z=y,\quad d_i z=1\quad(2\le i\le r+1),
 \qquad x,y\in Z_r,
\]
then $u=z(s_0y)^{-1}$ belongs to $N_{r+1}$ and
$d_0u=xy^{-1}$. Conversely $u\in N_{r+1}$ produces the same
homotopy from $u s_0y$ whenever $d_0u=xy^{-1}$.
For $r=0$ this is the corresponding edge calculation; the Kan
condition composes edge paths. The multiplication on these
homotopy classes is induced by the degreewise group product,
by Proposition~17.2. Hence there is an isomorphism of groups
\[
 \pi_r|G|\cong Z_r/d_0N_{r+1}\quad(r\ge0).
                                                        \tag{A.131.10}
\]
The realization comparison in Theorem~16.6 identifies the simplicial
and ordinary homotopy groups used here.

The promised presentation is completely specified as follows:
\[
 \Pi_r(G)=
 \left\langle x_a\ (a\in Z_r)\ \left|\
 x_1=1,\quad x_ax_b=x_{ab}\ (a,b\in Z_r),\quad
 x_{d_0u}=1\ (u\in N_{r+1})\right.\right\rangle.
                                                        \tag{A.131.11}
\]
Enumerate the canonical binary codes, use the finite tests in
(A.131.8), and dovetail their pairs to enumerate the displayed
relators. The product and face indices are computed by the already
specified total algorithms. Thus no word in (A.131.11) is indexed
by an unspecified isotopy or an undecidable predicate.
The first two families of relators present $Z_r$ via its full
multiplication table: mapping $x_a$ to $a$ and $a$ to $x_a$ gives
inverse homomorphisms. The last family takes the quotient by the
normal subgroup $d_0N_{r+1}$, which is precisely the subgroup in
(A.131.10). This proves that every relation, and only the required
relations, is included.

For $r\ge1$ the quotient is abelian, as also proved in the first
Proposition~17.3 on p.~68 of the same source. Therefore one may
instead enumerate the free abelian generators $e_a$ and rows
\[
 e_1,\qquad e_a+e_b-e_{ab},\qquad e_{d_0u}
                                                        \tag{A.131.12}
\]
of a countable integer relation matrix. Abelianization adds no
further relation to the already abelian group (A.131.11).
The subgroup $d_0N_{r+1}$ is recursively enumerable, since its
preimages have a decidable membership test, but its membership is
not claimed decidable. In particular none of this asserts an
algorithm that answers every isotopy question in finite time.

Finally A.113 identifies $\Pi_0^{\mathrm{PL}}(K)$ with the
definable component group in every fixed o-minimal expansion.
In the stated dimension-at-least-five range the same result gives
a PL--TOP weak equivalence over the real field, and hence the full
TOP model and all homotopy presentations asserted above.
This uses the comparison with the singular complex of the TOP
group, not a literal real-points map from an arbitrary field.

\emph{Finite verification and the exact meaning of the output.}
The attached independent program checks a rational triangulated
subclass of (A.131.4) for $B$ a point and $q=0,1,2$ using exact
rational arithmetic. It constructs the common subdivisions in
(A.131.5) and affine preimages in (A.131.6), tests both projection
triangulations, and checks nonconstant jointly PL families
\[
 h_t(x)=x+\tfrac12\min(x,1-x,t_0,\ldots,t_q)
 \quad(q=1,2).
\]
It checks inverse identities, subdivision equality, all faces of the
two-simplex example, and degeneracies of a nonidentity zero-simplex;
it rejects five defective graph codes. All thirty checks passed,
and an independent author rerun passed in $0.842377$ seconds.
A second author run, saved separately, passed in $0.761886$ seconds,
and the root reviewer independently passed all thirty in $0.778375$
seconds.
The program does not claim to implement general polynomial
quantifier elimination or certify the infinite Moore theorem.

The result is an algorithm from a finite complex to a countable
simplicial group with decidable simplices and computable structure,
and thence to complete recursively enumerable group presentations.
The real and arbitrary-field concordance spaces are identified with
that model by the proved coefficient comparisons. This gives more
than a name for a concordance group: every finite candidate and
every displayed relation is checked by terminating algebraic
procedures. It remains an exhaustive presentation, not an evaluation
into finite classical invariants. This theorem alone did not identify
the computable map $G^{\mathrm{PL}}\to G^{\mathrm{SA}}$ as a weak
equivalence. That comparison, and comparison with every fixed
o-minimal expansion, are now proved in A.160. The precise effective
whole-space consequence, requiring no decision oracle for the
expansion, is stated and proved in A.161.
\end{proof}



\begin{saresult}{A.132}{Refuted}{The unrestricted real PL--TOP concordance comparison fails on the three-sphere}
Give the topological homeomorphism groups the compact-open topology,
and write $P_{\mathrm{TOP}}(S^3)$ for the singular complex of
\[
 \operatorname{Homeo}(S^3\times I\ \mathrm{rel}\ S^3\times\{0\}).
\]
The upper boundary is free. Write $P_{\mathrm{PL}}(S^3)$ for the
simplicial group of jointly finite PL concordances with the same
boundary convention. Then there is a weak equivalence of based spaces
\[
 P_{\mathrm{TOP}}(S^3)\simeq
 \Omega\bigl(\mathrm{TOP}(4)/\mathrm O(4)\bigr),
                                                        \tag{A.132.1}
\]
where $\mathrm{TOP}(4)=\operatorname{Homeo}(\mathbb R^4;0)$ is the
actual group of global homeomorphisms, and the quotient on the
right denotes $\operatorname{hofib}(B\mathrm O(4)\to
B\mathrm{TOP}(4))$. In particular,
\[
 \pi_2 P_{\mathrm{PL}}(S^3)=0,
 \qquad
 \pi_2 P_{\mathrm{TOP}}(S^3)=\mathbb Z/2.
                                                        \tag{A.132.2}
\]
Thus the natural real PL-to-TOP comparison is not a weak equivalence
in unrestricted dimension. Both groups have trivial $\pi_0$ and
$\pi_1$ in this example, so this discrepancy is genuinely a
higher-parameter discrepancy.
\end{saresult}
\begin{proof}
We prove the topological fibre model directly, including its relative
parameter convention. This avoids extending a PL germ theorem to TOP
without checking its hypotheses.

\smallskip\noindent
\textbf{The germ simplicial group and its two restriction maps.}
Let $D=D^4\subset\mathbb R^4$ be the closed unit ball, and put
\[
 A=\operatorname{Sing}\operatorname{Homeo}(D;0),\qquad
 T=\operatorname{Sing}\operatorname{Homeo}(\mathbb R^4;0).
\]
When $\mathrm O(4)$ or $\mathrm{SO}(4)$ occurs in a simplicial map
below, its singular complex is understood. Weak equivalences of
the resulting simplicial sets are equivalently weak equivalences
of their realizations.
Let $G$ be the following simplicial group of parameterized germs.
A $q$-simplex is the germ along $\{0\}\times\Delta^q$ of a
homeomorphism between open neighborhoods of that set in
$\mathbb R^4\times\Delta^q$, preserving the parameter projection
and fixing the zero section. Neighborhoods and openness are relative
to $\mathbb R^4\times\Delta^q$. Two representatives are equivalent
when they agree on a neighborhood of the whole zero section.
Composition, inverse, and the simplicial operators are defined on
germs. Compactness of $\Delta^q$ gives a common positive source
radius and a common positive image radius for each representative;
its inverse is jointly continuous on the latter neighborhood.
There are germ-restriction maps
\[
 r:A\longrightarrow G,\qquad g:T\longrightarrow G.
                                                        \tag{A.132.3}
\]

\smallskip\noindent
\textbf{A continuous local extension operator.}
We use the following precise consequence of
\cite[Theorem~5.1, p.~70; proof of Corollary~1.2, p.~79]
{EdwardsKirby1971}. Given concentric balls
$C\subset\operatorname{int}U\subset U\subset\operatorname{int}D$,
where $C$ and $U$ are compact, there is a neighborhood $\mathcal V$
of the inclusion in the space of embeddings $U\to\operatorname{int}D$
and a continuous map
\[
 \sigma:\mathcal V\longrightarrow
 \operatorname{Homeo}(D\ \mathrm{rel}\ \partial D)
\]
such that $\sigma(v)|_C=v|_C$ and $\sigma(\mathrm{id})=\mathrm{id}$.
The extensions may be chosen supported in a fixed compact subset of
$\operatorname{int}D$ after shrinking $\mathcal V$.
Indeed, Theorem~5.1 supplies a deformation, continuous in the input
embedding and fixed at the inclusion, to an embedding equal to the
identity on $C$ and unchanged near $\partial U$. Composing the input
with the inverse of this final embedding, and extending by the
identity, is exactly the local ambient extension in the proof on
p.~79. Its support stays inside the image of $U$; embeddings
sufficiently close to the inclusion have these images in a fixed
compact subset of $\operatorname{int}D$.
This is a continuous operator on an open set of embeddings, not
merely a separate extension selected for each parameter.

\smallskip\noindent
\textbf{Relative lifting over a horn, with uniform radii.}
We show that both maps in (A.132.3) are Kan fibrations. Let
$H=\Lambda_i^q\subset\Delta^q$, $q\geq1$, and suppose a horn
lift $a:H\to\operatorname{Homeo}(D;0)$ and a germ simplex $f$
are given with $r(a)=f|_H$. The proof for global homeomorphisms
is identical. Choose a representative $f_t$ on a common neighborhood
of $0$ for all $t\in\Delta^q$. Equality of germs on each of the
finitely many horn faces, followed by one common shrinking, gives
\[
 f_t=a_t\quad\hbox{on a fixed ball about $0$, for every }t\in H.
                                                        \tag{A.132.4}
\]
Choose a strong deformation $d_s:\Delta^q\to\Delta^q$, $0\leq s\leq1$,
with $d_0(\Delta^q)\subset H$, $d_1=\mathrm{id}$, and
$d_s|_H=\mathrm{id}$. Such a deformation exists because a simplex
collapses onto any horn. Set
\[
 a^0_t=a_{d_0(t)},\qquad
 k_{t,s}=f_{d_s(t)}\circ(a^0_t)^{-1}
\]
as germs. Joint continuity of $a^0$ and its inverse, compactness of
$\Delta^q\times I$, and (A.132.4) permit these germs to be
represented by one family of local homeomorphisms on a fixed ball,
with their inverses defined on another fixed ball. They satisfy,
on a common possibly smaller ball,
\[
 k_{t,0}=\mathrm{id}\quad(t\in\Delta^q),\qquad
 k_{t,s}=\mathrm{id}\quad(t\in H,\ s\in I).
                                                        \tag{A.132.5}
\]
Here the second equality is actual equality on the common ball:
$d_s(t)=d_0(t)=t$ on $H$ and (A.132.4) applies. It is not only
pointwise equality of unparameterized germs.

Choose a compact ball $U\subset\operatorname{int}D$ sufficiently
small that every $k_{t,s}^{-1}$ is defined on $U$ and its image
lies in the common source domain of every $k_{t,v}$.
This is possible by compactness and the joint inverse continuity
near the zero section. Shrink $U$ also so that (A.132.5) holds
on all relevant source and image neighborhoods. Choose a smaller
concentric compact ball $C\subset\operatorname{int}U$, and then
a still smaller ball $V$ such that
\[
 k_{t,s}(V)\subset\operatorname{int}C
 \quad\hbox{for all }(t,s)\in\Delta^q\times I.
                                                        \tag{A.132.6}
\]
The last choice follows from joint continuity and the fact that
every map fixes $0$.

On the fixed domain $U$, the increments
\[
 v_{t;s,s'}=k_{t,s}\circ k_{t,s'}^{-1}|_U
\]
depend continuously on $(t,s,s')$ in the embedding space wherever
they are used, and equal the inclusion when $s=s'$.
For the single neighborhood $\mathcal V$ in the local extension
operator above, compactness of the diagonal
$\{(t,s,s):t\in\Delta^q,s\in I\}$ gives $\eta>0$ such that
\[
 |s-s'|<\eta\quad\Longrightarrow\quad
 v_{t;s,s'}\in\mathcal V
 \quad\hbox{for all }t.
                                                        \tag{A.132.7}
\]
In particular these increments have images in $\operatorname{int}D$.
This supplies one time subdivision for all auxiliary parameters;
there is no parameter-dependent choice of subdivision or extension.

Choose $0=s_0<\cdots<s_N=1$ with mesh less than $\eta$.
Starting with $b_{t,0}=\mathrm{id}$, define for
$s_j\leq s\leq s_{j+1}$
\[
 b_{t,s}=
 \sigma(v_{t;s,s_j})\circ b_{t,s_j}.
                                                        \tag{A.132.8}
\]
The formulas agree at the subdivision points because
$\sigma(\mathrm{id})=\mathrm{id}$, so this is a continuous
family of disk homeomorphisms. By induction, (A.132.6) and the
extension property on $C$ give
\[
 b_{t,s}|_V=k_{t,s}|_V.
\]
Indeed, if $x\in V$, the input to the new extension is
$b_{t,s_j}(x)=k_{t,s_j}(x)\in C$.
The extensions fix $0$, since the increments do and their
extensions agree on $C$ containing $0$. For $t\in H$ every
increment in (A.132.8) is the identity by (A.132.5), hence
$b_{t,s}=\mathrm{id}$ for all $s$. Consequently
\[
 \widetilde a_t=b_{t,1}\circ a^0_t
\]
is a continuous lift of the given germ simplex and agrees exactly
with $a_t$ on $H$. Its germ is $f_t$ on a common ball, by one
further uniform shrinking under the family $a^0$ if necessary.
This proves the required horn lifting. For $T$, extend each
$b_{t,s}$ by the identity outside $D$ and use the same formula.
Thus arbitrary germ paths and compact auxiliary parameter families
are treated by finitely many small increments; no assumption that
the original germ family is close to the identity is made.

\smallskip\noindent
\textbf{The disk germ fibre has a free outer face.}
Let $K=\ker r$. For $0<\epsilon<1$, put
\[
 K_\epsilon=
 \operatorname{Sing}\operatorname{Homeo}(D\ \mathrm{rel}\ D_\epsilon),
 \qquad D_\epsilon=\epsilon D.
\]
Then $K$ is the filtered union of the $K_\epsilon$, because a
simplex in the identity germ fibre is the identity on a common
neighborhood of $0$. Restriction to the annulus identifies
$K_\epsilon$ with $P_{\mathrm{TOP}}(S^3)$: the inner sphere is
fixed pointwise, and the outer sphere is unrestricted. Extension
by the identity on $D_\epsilon$ is its continuous inverse.

For $0<\epsilon<\epsilon_0<1$, the inclusion
$K_{\epsilon_0}\to K_\epsilon$ is homotopic to an isomorphism.
Choose a radial disk homeomorphism $\rho$ taking
$D_{\epsilon_0}$ onto $D_\epsilon$ and a path $\rho_u$ from the
identity to $\rho$ for which
$\rho_u(D_{\epsilon_0})=D_{v(u)}$ with
$\epsilon\leq v(u)\leq\epsilon_0$. For example, take increasing
piecewise affine radius functions with this prescribed value at
$\epsilon_0$ and linearly interpolate the radius functions.
The homotopy $h\mapsto\rho_u h\rho_u^{-1}$ takes values in
$K_\epsilon$, starts at the inclusion, and ends at conjugation
by $\rho$, an isomorphism onto $K_\epsilon$.
Every finite simplicial sphere and every finite simplicial
nullhomotopy in $K$ uses finitely many simplices and therefore
one common radius. Filtered unions of these Kan complexes thus
give
\[
 K_{\epsilon_0}\xrightarrow{\simeq}K,
 \qquad K\simeq P_{\mathrm{TOP}}(S^3).
                                                        \tag{A.132.9}
\]
This identification imposes neither a fixed outer boundary nor
an additional factor recording a boundary parametrization.

\smallskip\noindent
\textbf{Global germs, with the component issue made explicit.}
The strict fibre $L=\ker g$ is weakly contractible. In fact, for
a compact parameter family $h_t$ in this fibre, all $h_t$ are the
identity on some $D_\epsilon$. The formula
\[
 c_s(h_t)(x)=
 \begin{cases}s^{-1}h_t(sx),&s>0,\\x,&s=0,\end{cases}
                                                        \tag{A.132.10}
\]
contracts the family to the identity in the compact-open topology.
On every fixed compact subset of $\mathbb R^4$, it is exactly the
identity for all sufficiently small $s$, uniformly in $t$; the
same assertion holds for inverses. For $0<s\leq1$ it is the
identity on $D_\epsilon$, and hence remains a family in $L$.
The formula fixes any parameters at which the input is the identity.

Let $T_0$ and $G_0$ be the identity components as simplicial sets.
The Kan fibration $g$ restricts to a Kan fibration
$T_0\to G_0$: a lift whose horn is in the identity component
remains there, because a simplex is connected. Every vertex of
$G_0$ can be reached by finitely many edges from the identity,
and the $q=1$ lifting just proved supplies a lift in $T_0$.
Together with (A.132.10), the fibre exact sequence proves
\[
 T_0\xrightarrow{\simeq}G_0.
                                                        \tag{A.132.11}
\]
No surjectivity assertion about all components of the germ group,
and no theorem globalizing an arbitrary isolated germ, is required.

\smallskip\noindent
\textbf{The disk group and the actual topological three-sphere group.}
Boundary restriction
$\operatorname{Homeo}(D;0)\to\operatorname{Homeo}(S^3)$ has the
radial-coning section. Its kernel is contracted by the relative
Alexander isotopy
\[
 a_u(x)=
 \begin{cases}u\,a(x/u),&\|x\|\leq u,\ u>0,\\
 x,&\|x\|\geq u\text{ or }u=0.
 \end{cases}
\]
These formulas agree on $\|x\|=u$ because $a$ fixes $\partial D$;
they fix $0$, and continuity at $u=0$ is uniform, with displacement
at most $2u$. The same bound holds for inverses. Splitting a
homeomorphism by its coned boundary restriction therefore proves
that boundary restriction is a homotopy equivalence.

The natural inclusion
$\mathrm O(4)\to\operatorname{Homeo}(S^3)$ is a homotopy
equivalence by \cite[Appendix, statement~(16), p.~606]{Hatcher1983}.
The original explicitly states the TOP version there; the
comparison discussion at statement~(3), p.~605 also explicitly
permits TOP in place of PL. Thus this is a statement about the
actual $C^0$ topological group, not an inference from the smooth
Smale theorem alone and not a statement about slice-wise PL maps.
It follows that
\[
 \mathrm O(4)\xrightarrow{\simeq}A,
 \qquad \mathrm{SO}(4)\xrightarrow{\simeq}A_+,
                                                        \tag{A.132.12}
\]
where $A_+$ is the orientation-preserving, hence identity,
component. An element of $K$ fixes a neighborhood of $0$ and is
orientation preserving, so the fibre of $r:A_+\to G_0$ is still
all of $K$.

Extend $a\in\operatorname{Homeo}(D;0)$ to a global homeomorphism
$e(a)$ by setting, outside $D$,
\[
 e(a)(r\omega)=r\,a|_{S^3}(\omega)\quad(r\geq1).
\]
This is a continuous group map, satisfies $g e=r$, and on
$\mathrm O(4)$ is the usual linear action on $\mathbb R^4$.
By (A.132.11), (A.132.12), and the fibre property already proved,
\[
 K\simeq\operatorname{hofib}(A_+\to T_0)
 \simeq\operatorname{hofib}(\mathrm{SO}(4)\to T_0)
 =\operatorname{hofib}_{1}(\mathrm O(4)\to T).
                                                        \tag{A.132.13}
\]
For the last equality, an orientation-reversing orthogonal map
cannot be joined to the identity in $T$. Every orientation-preserving
orthogonal map is in its identity component. Only those components
enter the homotopy fibre over $1$; no assertion about other
components of $T$ is needed.
For simplicial groups the contractible universal bundle identifies
$\Omega BG$ with $G$ up to weak equivalence. Applying this to the
homotopy fibre of $B\mathrm O(4)\to BT$ identifies the last space
in (A.132.13) with
$\Omega(\mathrm{TOP}(4)/\mathrm O(4))$. This and (A.132.9) prove
(A.132.1).

\smallskip\noindent
\textbf{The two homotopy calculations.}
The stabilization
\[
 \mathrm{TOP}(4)/\mathrm O(4)\longrightarrow\mathrm{TOP}/\mathrm O
\]
is $5$-connected by \cite[Theorem~8.7A, p.~128]{FreedmanQuinn1990}.
That source writes its comparison with $\mathrm{DIFF}(4)$;
for the precise global-group and orthogonal-group convention here,
we use the corrected modern account
\cite[\S1.1, Theorem~1.3 and the immediately following paragraph]
{GabaiGayHartmanKrushkalPowell2026} states the same
$5$-connected stabilization using exactly the global group
$\mathrm{TOP}(4)=\operatorname{Homeo}(\mathbb R^4;0)$ and the
classifying-space homotopy fibre convention used here.
It supplies the correction to Quinn's argument needed for the
full original theorem; the unchecked replacement criterion is
not used as an input.
By \cite[Essay~V, Theorem~5.5, p.~251]{KirbySiebenmann1977},
\[
 \pi_i(\mathrm{TOP}/\mathrm O)=
 \begin{cases}\mathbb Z/2,&i=3,\\0,&0\leq i\leq6,\ i\ne3.
 \end{cases}
\]
In particular (A.132.1) gives
$\pi_2P_{\mathrm{TOP}}(S^3)=\mathbb Z/2$ and vanishing
$\pi_0,\pi_1$.

For PL, set $X_j=\mathrm{PL}_j/\mathrm O_j$ as in A.128.
Its checked, jointly PL, free-top sphere fibre model is
\[
 P_{\mathrm{PL}}(S^3)\simeq
 \operatorname{hofib}(X_3\longrightarrow X_4).
\]
The low-degree formula
\cite[\S5, p.~37, formula~(3)]{BurgheleaLashof1974II} gives
$\pi_iX_j\cong\pi_i(\mathrm{PL}/\mathrm O)$ for $i<j$.
The latter groups vanish through degree six
\cite[Annex~C, p.~318]{KirbySiebenmann1977}.
The fibre exact sequence therefore has zero outer terms in
\[
 \pi_3X_4\longrightarrow\pi_2P_{\mathrm{PL}}(S^3)
 \longrightarrow\pi_2X_3,
\]
proving its middle group zero. The same argument in degrees
zero and one gives the asserted lower vanishing. The natural
PL-to-TOP inclusion consequently cannot induce an isomorphism
on $\pi_2$, proving the refutation.

For an explicit finite base, take the boundary of the abstract
$4$-simplex. The accompanying exact certificate checks its
$f$-vector $(5,10,10,5)$, every face link, integral boundary
Smith forms, and a primitive fundamental cycle: $53$ checks
pass. This computation certifies the finite PL three-sphere;
the existence and detection of the non-PL topological family
are proved by the fibre calculation and the cited stabilization
theorem, not asserted to have been numerically tested.
\end{proof}
